% !TEX root = /home/eb/.vscode/extensions/qftrules.programme-colle/tmp/Exercice.tex
%%%%%%%%%%%%%%%%%%%%%%%%%%%%%%%%%%%%%%%%%%%%%
\begin{exo}[PCSI][1][colle]{Particule dans des champs E et B orthogonaux}
  On étudie dans le référentiel terrestre supposé galiléen le mouvement d'une particule chargée de charge $q$ et de masse $m$ initialement immobile et à l'origine du repère, dans un champ électromagnétique uniforme et stationnaire $\vv{E}=E\,\vv{e_y}$ et $\vv{B}=B\,\vv{e_z}$.
  \begin{quest}
  \item Déterminer les équations différentielles régissant le mouvement de la particule. En déduire les équations horaires de la trajectoire.
  \item En étudiant certains instants bien choisis, déterminer la nature de la trajectoire de la particule.
  \end{quest}
  \end{exo}
  %%%%%%%%%%%%%%%%%%%%%%%%%%%%%%%%%%%%%%%%%%%%%%%%%%%%%%%%%%%%%%%%%%%%%%%%%%%%%%%%

%%%%%%%%%%%%%%%%%%%%%%%%%%%%%%%%%%%
\begin{exo}[PC][1][devoir]{Étude du mouvement d'un satellite}
  \Annale[Physique MP]{CCINP}{2025}

  \begin{wrapfigure}{r}{0.42\columnwidth}
  \includegraphics[width=0.42\columnwidth]{2025_09_19_1f7500c51cd45be5f039g-1(1)}
\end{wrapfigure}

Les systèmes d'observation des océans par satellite ont été imaginés et développés au début des années 70. Depuis, plus d'une quinzaine de satellites d'observation embarquant des altimètres radars ont été lancés dans le but d'observer le comportement des océans.

Issues d'une coopération du CNES et de la NASA, la série des satellites Topex-Poséidon, initiée en 1992, puis celle des satellites Jason, ont permis de mesurer l'élévation moyenne des mers avec précision : $(3,6 \pm 0,1) \mathrm{mm} / \mathrm{an}$ durant ces trente dernières années.

% \begin{figure}[h]
% \begin{center}

% \captionsetup{labelformat=empty}
% \caption{Figure 2 - Satellites altimétriques lancés depuis 1992. Vue d'artiste. Crédit : CNES}
% \end{center}
% \end{figure}

On se propose dans cette partie d'étudier le mouvement d'un tel satellite, en orbite autour du centre $\rm O$ de la Terre, modélisée par un corps de répartition de masse à symétrie sphérique, de rayon $R_{T}$ et de masse $M_{T}$.

\subsection{Force centrale conservative}
On commence par étudier le mouvement d'un mobile quelconque, de masse $m$ et assimilé à un point matériel $\rm M$, dans le référentiel géocentrique ($\mc{R}_{T}$) considéré comme galiléen. Le mobile n'est soumis qu'à la seule action de la Terre.
\begin{quest}
  \QR{   Rappeler la définition du référentiel géocentrique et celle d'un référentiel galiléen.
  }{
    Le référentiel géocentrique est le référentiel lié au centre de la Terre.
Un référentiel galiléen est un référentiel dans lequel le principe d'inertie est vérifié.
  }

  \QR{ Rappeler l’expression du champ de gravitation terrestre $\Vec{\mc{G}}(M)$ en un point $M$ extérieur à la Terre en fonction de la constante de gravitation universelle $G$, de la masse $M_{T}$, de la distance $r=O M$ et du vecteur unitaire $\Vec{u}=\Vec{O M}/r$. En déduire l'expression $\Vec{F}_{g}$ de la force de gravitation exercée par la Terre sur le mobile de masse $m$.
  }{
   Le champ gravitationnel s'exprime donc : 
   \eqbox{
   \Vec{\mc{G}}(r)=-\frac{G M_{T}}{r^{2}} \Vec{u}.
   }
   La force gravitationnelle s'exprime alors : 
   \eqbox{
\Vec{F}_{g}=m \Vec{\mc{G}}=-\frac{G m M_{T}}{r^{2}} \Vec{u}.
   }
  }

  \QR{   Montrer que le moment cinétique $\Vec{\mc{L}_{O}}$ du mobile par rapport au point $\rm O$ est une constante du mouvement. En déduire que la trajectoire du mobile est plane.
  }{
   Le mobile n'est soumis qu'à la force gravitationnelle $\Vec{F}_{g}$.
Le théorème du moment cinétique par rapport à $O$ s'écrit alors : $$\frac{d \Vec{\mc{L}}_{O}}{d t}=\Vec{\mc{M}}_{O}\left(\Vec{F}_{g}\right).$$
Or on a : $$\Vec{\mc{M}}_{O}\left(\Vec{F}_{g}\right)=\Vec{O M} \wedge\left(-\frac{G m M_{T}}{r^{2}} \Vec{u}\right)=(r \Vec{u}) \wedge\left(-\frac{G m M_{T}}{r^{2}} \Vec{u}\right)=\Vec{0}.$$
On a donc: $\frac{d \Vec{\mc{L}}_{O}}{d t}=\Vec{0}$ et donc : 
\eqbox{
\Vec{\mc{L}}_{O}=\Vec{\text{cte}}.
}
Par définition du moment cinétique, on a : $\Vec{\mc{L}}_{O}=\Vec{O M} \wedge m \Vec{v}$.
Le moment cinétique est constant et orthogonal à la fois à $\Vec{O M}$ et à $\Vec{v}$. On a donc nécessairement le point $M$ et le vecteur $\Vec{v}$ inclus dans le plan orthogonal à $\Vec{\mc{L}}_{O}$ passant par $O$. La trajectoire du mobile est donc \fbox{plane}.
  }
\end{quest}

Dans la suite, on associera au référentiel ($\mc{R}_{T}$) le repère orthonormé ($O, \Vec{e_{x}}, \Vec{e_{y}}, \Vec{e_{z}}$) de façon à ce que le moment cinétique $\Vec{\mc{L}_{O}}$ soit aligné avec $\Vec{e_{z}}$. On posera $\Vec{\mc{L}_{O}}=\mc{L}_{0} \Vec{e_{z}}$ et on se placera en coordonnées polaires ( $r, \theta$ ), de centre $O$, pour décrire le mouvement du mobile.
% \begin{figure}[h]
% \begin{center}
\begin{fig}[label]
  \centering
  \includegraphics[width=0.5\textwidth]{2025_09_19_1f7500c51cd45be5f039g-1}
  %\capt{}
\end{fig}
% \captionsetup{labelformat=empty}
% \caption{Figure 3 - Description du mouvement du mobile dans le système de coordonnées polaires}
% \end{center}

\begin{quest}
  \QR{   Montrer que la force gravitationnelle s'exerçant sur le mobile dérive d'une énergie potentielle $\mc{E}_{p}$. Établir l'expression de celle-ci en la prenant, par convention, nulle à l'infini.
  }{
 La force gravitationnelle dans les coordonnées polaires considérées s'écrit : $$\vec{F}_{g}=-\frac{G m M_{T}}{r^{2}} \vec{e}_{r}.$$
La force gravitationnelle est une force conservative, on a donc : $$\vec{F}_{g}=-\Vec{\operatorname{grad}}\left(\mc{E}_{p}\right).$$
Par projection selon l'axe radial, on a : $$-\frac{G m M_{T}}{r^{2}}=-\frac{d \mc{E}_{p}}{d r}.$$
D'où : $\mc{E}_{p}=+\int \frac{G m M_{T}}{r^{2}} d r=-\frac{G m M_{T}}{r}+$ cte.
En prenant l'énergie potentielle nulle à l'infini, on a finalement : 
\eqbox{
\mc{E}_{p}=-\frac{G m M_{T}}{r}.
}
  }
  % \end{figure}
  
  \QR{   Montrer que l'énergie mécanique $\mc{E}_{m}$ est une constante du mouvement et qu'elle peut se mettre sous la forme :
  \begin{equation*}
  \mc{E}_{m}=\frac{1}{2} m \dot{r}^{2}+\mc{E}_{p, \text{eff}}(r) \tag{1}
  \end{equation*}
  où $\mc{E}_{p, \text {eff}}(r)$ est un terme, appelé énergie potentielle effective, que l'on exprimera en fonction de $G, m, M_{T}, \mc{L}_{0}$ et de $r$.
  }{
La seule force qui s'exerce sur le mobile est une force conservative, d'après le théorème de l'énergie mécanique, l'énergie mécanique du mobile se conserve.

L'énergie mécanique s'écrit : $$\mc{E}_{m}=\mc{E}_{c}+\mc{E}_{p}=\frac{1}{2} m v^{2}-\frac{G m M_{T}}{r}=\frac{1}{2} m\left(\dot{r}^{2}+(r \dot{\theta})^{2}\right)-\frac{G m M_{T}}{r}.$$
Le moment cinétique s'écrit : $$\Vec{\mc{L}}_{O}=\Vec{O M} \wedge m \vec{v}=m r^{2} \dot{\theta} \vec{e}_{Z}.$$ D'où : $$\dot{\theta}=\frac{\mc{L}_{O}}{m r^{2}}.$$
On a donc : $$\mc{E}_{m}=\frac{1}{2} m \dot{r}^{2}+\frac{1}{2} m r^{2} \dot{\theta}^{2}-\frac{G m M_{T}}{r},$$
soit 
\eqbox{
\mc{E}_{m} =\frac{1}{2} m \dot{r}^{2}+\frac{\mc{L}_{O}^{2}}{2 m r^{2}}-\frac{G m M_{T}}{r}.
}
}


\QR{   Expliquer pourquoi l'énergie mécanique du mobile est nécessairement supérieure ou égale à son énergie potentielle effective.
}{
Dans l'expression de l'énergie mécanique, le terme $\frac{1}{2} m \dot{r}^{2}$ est positif ou nul.
On a donc bien : $\mc{E}_{m}=\frac{1}{2} m \dot{r}^{2}+\mc{E}_{p, \text { eff }}(r) \geq \mc{E}_{p, \text { eff }}(r)$.
Il y a deux positions pour lesquelles $\mc{E}_{m}=\mc{E}_{p, \text { eff }}$.
}

\QR{   Représenter graphiquement, pour une valeur donnée de $\mc{L}_{0}$, l'énergie potentielle effective $\mc{E}_{p, \text {eff}}$ du mobile en fonction de $r$. Faire apparaître sur le graphique l'énergie mécanique d'une trajectoire associée à un état lié. On rappelle que, pour une force centrale en $1 / r^{2}$, la trajectoire d'un état lié est elliptique.
}{
On réalise le graphique demandé.\\
On trace la forme habituelle d'un tel potentiel effectif.\\
Un état lié correspond à un état où le rayon $r$ est borné, soit ici un cas où $\mc{E}_{m}<0$.

\begin{fig}[label]
  \centering
  \includegraphics[width=0.7\textwidth]{2025_09_23_8476f4d4293d04e17198g-02}
  %\capt{}
\end{fig}
}

\QR{   Pour un mouvement elliptique quelconque, indiquer à quelles positions particulières l'énergie mécanique est égale à l'énergie potentielle effective. Caractériser le mouvement du mobile dans le cas où l'énergie mécanique est égale au minimum de l'énergie potentielle effective.
}{
Il y a deux positions pour lesquelles $\mc{E}_{m}=\mc{E}_{p, \text {eff}}$.
\begin{itemize}
  \item 
  
  La plus proche, de rayon $r_{P}$, est le périgée. Le périgée correspond au point le plus proche du centre attracteur d'une orbite elliptique.
  
  \item La plus éloignée, de rayon $r_{A}$, est l'apogée. L'apogée correspond au point le plus éloigné du centre attracteur d'une orbite elliptique.
\end{itemize}

Dans le cas où l'énergie mécanique est égale à l'énergie potentielle effective minimale, le rayon $r=r_{0}$ est fixé. On a alors un mouvement \fbox{circulaire}.
}

\end{quest}
La plupart des mesures effectuées par les satellites altimétriques se font à partir de l'orbite altimétrique de référence, que l'on considérera ici comme une orbite circulaire de rayon $R$. Dans la suite, le mobile étudié correspond à un satellite altimétrique de masse $m$, assimilable à un point matériel.
\begin{quest}

\QR{   Exprimer l'énergie mécanique $\mc{E}_{m, \text{alt}}$ du satellite situé sur l'orbite altimétrique de référence, en fonction de $G, M_{T}, m$ et de $R$.
}{
On considère un mouvement circulaire de rayon $R$.
Le principe fondamental de la dynamique s'écrit : $m \vec{a}=\vec{F}_{g}$.
On a alors: $$m\left(-R \dot{\theta}^{2} \vec{e}_{r}+R \ddot{\theta} \vec{e}_{\theta}\right)=m\left(-\frac{v^{2}}{R} \vec{e}_{r}+\dot{v} \vec{e}_{\theta}\right)=-\frac{G m M_{T}}{R^{2}} \vec{e}_{r}.$$
Les projections permettent d'écrire: $-m \frac{v^{2}}{R}=-\frac{G m M_{T}}{R^{2}}$ et $\dot{v}=0$.

D'où : $v=\sqrt{\frac{G M_{T}}{R}}=$ cte.
L'énergie mécanique s'écrit alors : $$\mc{E}_{m, \text { alt }}=\mc{E}_{c}+\mc{E}_{p}=\frac{1}{2} m v^{2}-\frac{G m M_{T}}{R}=\frac{G m M_{T}}{2 R}-\frac{G m M_{T}}{R},$$
soit 
\eqbox{
\mc{E}_{m, \text { alt }}=-\frac{G m M_{T}}{2 R}.
}
}

\QR{  Établir la troisième loi de Kepler dans le cas particulier d'une orbite circulaire, en utilisant les paramètres liés à l'orbite altimétrique.
}{
La période de révolution du mobile s'écrit : $T=\frac{2 \pi R}{v}=2 \pi R \times \sqrt{\frac{R}{G M_{T}}}=\sqrt{\frac{4 \pi^{2} R^{3}}{G M_{T}}}$.\\
On retrouve alors l'expression de la troisième loi de Kepler pour une orbite circulaire : 

\eqbox{
\frac{T^{2}}{R^{3}}=\frac{4 \pi^{2}}{G M_{T}}.
}
}

\end{quest}
On admettra que la troisième loi de Kepler est valable plus généralement pour un mouvement elliptique. Son expression peut se déduire de l'équation obtenue pour le mouvement circulaire, en remplaçant le rayon $R$ de l'orbite circulaire par le demi-grand axe $a$ de la trajectoire elliptique.

\subsection{Jason-2 : un exemple pour la fin de vie des satellites}
En fin de vie, pour que ne soit pas laissé un objet non contrôlé sur l'orbite altimétrique de référence, le satellite Jason-2 a été dirigé vers une orbite dite «cimetière», d'altitude légèrement moins haute que celle de l'orbite altimétrique de référence, avant d'être définitivement abandonné. On se propose dans cette sous-partie d'étudier le cas d'une manœuvre de ce type dans le cas très simplifié, illustré figure 4, d'un transfert entre deux orbites circulaires coplanaires sous la seule action de l'attraction terrestre. L'orbite de transfert, appelée orbite de Hohmann, correspond à une ellipse dont l'un des foyers est le centre $\rm O$ de la Terre, dont l'apogée $\rm A$ est situé sur l'orbite altimétrique de référence (rayon $R$) et dont le périgée $\rm P$ est sur l'orbite cimetière (rayon $R_{c}$).

Pour modifier l'orbite du satellite, il faut l'accélérer ou le freiner en commandant le fonctionnement et la direction de ses moteurs. On considérera que la poussée générée par ceux-ci s'exerce pendant une durée tellement courte que les changements d'orbites se font instantanément.

% \begin{figure}[h]
% \begin{center}
\begin{fig}[label]
  \centering
  \includegraphics[width=0.5\textwidth]{2025_09_19_1f7500c51cd45be5f039g-3}
  %\capt{}
\end{fig}
% \captionsetup{labelformat=empty}
% \caption{Figure 4 - Tracé des différentes orbites du satellite}
% \end{center}
% \end{figure}
\begin{quest}
  \QR{  Montrer que l'énergie mécanique $\mc{E}_{m, \mathrm{tr}}$ du satellite sur l'orbite de transfert se met sous la forme :
  \begin{equation}
    \mc{E}_{m, \mathrm{tr}}=-\frac{G M_{T} m}{R+R_{c}}
  \end{equation}
  }{
  On part de l'expression de l'énergie mécanique trouvée plus haut : $$\mc{E}_{m}=\frac{1}{2} m \dot{r}^{2}+\mc{E}_{p, \text { eff }}(r).$$

Au niveau du périgée et de l'apogée la trajectoire est tangente à $\Vec{O M}$, et donc : $\dot{r}=0$.

On a alors : $$\mc{E}_{m}=\mc{E}_{p, \mathrm{eff}}\left(r_{P}\right)=\mc{E}_{p, \mathrm{eff}}\left(r_{A}\right).$$


Au périgée on peut donc écrire : $$\mc{E}_{m}=\frac{\mc{L}_{O}^{2}}{2 m R_{c}^{2}}-\frac{G m M_{T}}{R_{c}}.$$ D'où : $$\frac{\mc{L}_{O}^{2}}{2 m}=R_{c}^{2}\left[\mc{E}_{m}+\frac{G m M_{T}}{R_{c}}\right].$$


Et à l'apogée : $$\mc{E}_{m}=\frac{\mc{L}_{O}^{2}}{2 m R^{2}}-\frac{G m M_{T}}{R}.$$ D'où : $$\frac{\mc{L}_{O}^{2}}{2 m}=R^{2}\left[\mc{E}_{m}+\frac{G m M_{T}}{R}\right].$$

Le moment cinétique se conservant, on a : 
$$R_{c}^{2}\left[\mc{E}_{m}+\frac{G m M_{T}}{R_{c}}\right]=R^{2}\left[\mc{E}_{m}+\frac{G m M_{T}}{R}\right].$$

Et donc : $$\mc{E}_{m}\left(\frac{1}{R^{2}}-\frac{1}{R_{c}^{2}}\right)=\frac{G m M_{T}}{R R_{c}^{2}}-\frac{G m M_{T}}{R_{c} R^{2}}.$$


D'où finalement : 
$$\mc{E}_{m}=\left[\frac{G m M_{T}}{R R_{c}^{2}}-\frac{G m M_{T}}{R_{c} R^{2}}\right] \cdot \frac{R^{2} R_{c}^{2}}{R_{c}^{2}-R^{2}}=\frac{G m M_{T} R}{R_{c}^{2}-R^{2}}-\frac{G m M_{T} R_{c}}{R_{c}^{2}-R^{2}}=G m M_{T} \frac{R-R_{c}}{\left(R_{c}-R\right)\left(R+R_{c}\right)},$$
soit 
\eqbox{
\mc{E}_{m}=-\frac{G m M_{T}}{R+R_{c}}.
}
On retrouve bien l'expression demandée par le sujet.
  }

\QR{  Exprimer la variation d'énergie mécanique $\Delta \mc{E}_{m}=\mc{E}_{m, \mathrm{tr}}-\mc{E}_{m, \text {alt}}$ nécessaire pour passer de l'orbite initiale à l'orbite de transfert. Commenter le signe de $\Delta \mc{E}_{m}$.
}{
 La variation d'énergie mécanique s'écrit : $$\Delta \mc{E}_{m}=\mc{E}_{m, \mathrm{tr}}-\mc{E}_{m, \mathrm{alt}}=-\frac{G m M_{T}}{R+R_{c}}-\left(-\frac{G m M_{T}}{2 R}\right)=-G m M_{T} \frac{R-R_{c}}{2 R\left(R+R_{c}\right)}.$$
Le terme $R-R_{c}$ est positif, donc la variation d'énergie mécanique est négative : 
\eqbox{
\Delta \mc{E}_{m}<0.
}
Le système doit donc \fbox{perdre de l'énergie} pour passer sur l'orbite de transfert.
}

\QR{  En justifiant la réponse, indiquer s'il faut accélérer ou freiner le satellite pour le transférer en $\rm P$ de l'orbite de transfert à l'orbite cimetière.
}{
L'énergie mécanique sur l'orbite cimetière est : $$\mc{E}_{m, \mathrm{cim}}=-\frac{G m M_{T}}{2 R_{c}}.$$
La variation d'énergie mécanique entre l'orbite de transfert et l'orbite cimetière est donc : $$\Delta \mc{E}_{m}^{\prime}=\mc{E}_{m, \mathrm{cim}}-\mc{E}_{m, \mathrm{tr}}.$$
On a alors : $$\Delta \mc{E}_{m}^{\prime}=-\frac{G m M_{T}}{2 R_{C}}-\left(-\frac{G m M_{T}}{R+R_{C}}\right)=-G m M_{T} \frac{R-R_{C}}{2 R_{C}\left(R+R_{C}\right)}<0.$$
La variation d'énergie mécanique calculée est négative, le satellite doit donc perdre de l'énergie pour passer de l'orbite de transfert à l'orbite cimetière. Or au niveau du transfert le rayon est localement constant, le satellite doit donc perdre de l'énergie cinétique, soit être freiné en $P$.
}
\end{quest}

\end{exo}
%%%%%%%%%%%%%%%%%%%%%%%%%%%%%%%%%%%

%%%%%%%%%%%%%%%%%%%%%%%%%%%%%%%%%%%
\begin{exo}[PC][1][devoir]{Rayonnement gravitationnel d'un système de deux étoiles}
    Parmi les sources d'ondes gravitationnelles, l'effondrement d'un système binaire d'étoiles à neutrons
    est l'un des phénomènes que l'on pense détecter à l'aide de l'interféromètre décrit en première partie.
    Nous étudierons les aspects mécaniques de ce phénomène, dans le cadre simplifié de la dynamique
    Newtonienne.
    \subsection{Point matériel en rotation autour d'un astre}
    Le référentiel d'étude $\Rcal$, est supposé galiléen. Considérons dans cette question un astre $A$ de masse $M$, supposé immobile dans $\Rcal$, autour duquel
    gravite un petit objet $S$ de masse $m \ll M$. Désignons par $G$ la constante de gravitation universelle et
    notons $\Vec{r} = \Vec{AS}$ et $\er = \Vec{AS}/r$.
    \begin{quest}
    \item Qu'est ce qu'un référentiel galiléen ?

    \sol{
    Voir cours.
    }
    \item Rappeler l'expression de la force de gravitation subie par $S$ de la part de $A$. Montrer que cette
    force dérive d'une énergie potentielle $E_p(r)$ qui sera exprimée en prenant l'origine de l'énergie
    potentielle à l'infini.

    \sol{
    $\Vec{F} = -\frac{GMm}{r^2}\er$, de travail élémentaire $\delta W = \Vec{F} \cdot \Vec{\d OM} = -\frac{GMm}{r^2} \d r = -\d E_p$ avec $E_p(r) = -\frac{GMm}{r}$. On a choisi la constante d'intégration pour que $E_p(r=\infty)=0$.
    }

    \item Montrer que le moment cinétique de $S$ relativement au point $A$ dans $\Rcal$, noté $\Vec{L}_{A/\Rcal}$, est un
    vecteur constant. En déduire que la trajectoire de $S$ est située dans un plan passant par $A$.

    \sol{
    Le TMC donne $\dd{\Vec{L}_{A/\Rcal}}{t} = \Vec{M}(\Vec{F}) = \Vec{AM} \pv \Vec{F}=\Vec{0}$. On a $\Vec{L}_{A/\Rcal}=\Vec{AS}\pv m\Vec{v}(S) = \Vec{\cte}$, donc la trajectoire est contenue dans le plan orthogonal à $\Vec{L}_{A/\Rcal}$.
    }
    \item Admettons pour simplifier que le point $S$ ait une trajectoire circulaire de rayon $R$, avec
    une vitesse angulaire de rotation $\omega$. Exprimer $\omega$
    en fonction de $G$, $M$ et $R$.

    \sol{
    On a vu en TD que $v=r\omega=\sqrt{\frac{GM}{r}}$. On en déduit que $\omega = \sqrt{\frac{GM}{r^3}}$.
    }
    \item  Rappeler l'énoncé de la $3^{\rm e}$ loi de KEPLER relative à la période de rotation des satellites
    autour d'une étoile.
    Démontrez explicitement cette loi dans le cas d'une trajectoire circulaire.

    \sol{
    « Un satellite gravitant avec une période $T$ sur une orbite elliptique de demi-grand axe $a$ vérifie $\frac{T^2}{a^3}=\cte$. Pour une trajectoire circulaire, on a bien $T^2 = \frac{4\pi^2}{\omega^2} = 4\pi^2 \frac{r^3}{GM}$ et donc $\frac{T^2}{a^3} = \frac{4\pi^2}{GM}$.
    }

    \item Montrer que l'énergie mécanique du système s'écrit : $E_m = -GMm/2R$.
    Commenter le signe de cette énergie.

    \sol{
    On a $E_m =E c +E_p = \frac{1}{2}mv^2 -\frac{GMm}{r} = \frac{1}{2}\frac{GM}{r}-\frac{GMm}{r}$ d'où le résultat. On a $E_m<0$, ce qui traduit que le système astre-satellite est lié.
    }

    \end{quest}


    \subsection{Effondrement du système binaire}
    Considérons désormais l'ensemble formé par deux étoiles à neutron $A_1$ et $A_2$, de masses identiques $M={\SI{2.8e30}{kg}}$, distantes de $R$, en interaction
    gravitationnelle. Cet ensemble est supposé mécaniquement isolé. Peu de temps avant
    l'effondrement, elles ont une période de rotation très faible $T = \SI{0,1}{s}$. Le système binaire des deux étoiles $A_1$, $A_2$ est la source d'ondes gravitationnelles, qui transportent
    une certaine énergie. Un calcul de relativité générale montre que la puissance ainsi \guill{rayonnée} s'écrit, dans le référentiel $\Rcal$ :
    $P = KM^2R^4\omega^6$, où $K$ est une constante s'exprimant en fonction de $G$ et $c$ (vitesse de la lumière
    dans le vide) sous la forme : $K= 8G/5c^5$.  Du point de vue mécanique,
    l'émission des ondes gravitationnelles peut être modélisée par une force non conservative agissant
    sur le système des deux étoiles, avec une puissance $-P$.

    \begin{quest}
    \item D'après la section précédente, que vaut l'énergie mécanique $E_m$ ?
    \sol{
    On remplace $m$ par $M$, soit $E_m = -\frac{1}{2}\frac{GM^2}{R}$.
    }

    \item Qu'est ce qu'une force non conservative ? Quelle relation existe-t-il entre $\dd{E_m}{t}$ et $P$ ?
    Quelle est la conséquence de cette perte d'énergie sur la distance $R$ entre les deux étoiles ?

    \sol{
    Une force non conservative ne dérive pas d'une énergie potentielle (typiquement, il s'agit d'une force de frottements). Le théorème de la puissance mécanique donne
    $\dd{E_m}{t} = -P$. $E_m$ diminue et donc $R$ diminue : les astres se rapprochent.
    }

    \item Montrer que $R$ varie selon une loi : $\dd{R}{t} = -\frac{\alpha}{R^3}$ et exprimer $\alpha$ en fonction de $K$, $G$ et
    $M$.
    \sol{
    L'expression exacte (hors programme) de $\omega$ dans cette situation est $\omega = \sqrt{2GM/R^3}$ (on remplace $M$ par la somme des deux masses dans un système binaire). Mais j'acceptais (et j'attendais) l'expression trouvée au-dessus sans le facteur $2$.


    $\dd{E_m}{t} = -P$ donne $ -GM^2\dd{1/R}{t} =  \frac{GM^2}{R^2}\dd{R}{t} = - KM^2R^4(2GM/R^3)^3$ soit $\frac{G}{2R^2} \dd{R}{t} = -\frac{8KG^3M^3}{R^5}$ soit
    $ \dd{R}{t} = -\frac{\alpha}{R^3}$ avec $\alpha = 16 K G^2 M^3$.
    }

    \item Au temps $t = 0$, la distance entre les étoiles est $R(t = 0) = R_0=\SI{4.6e5}{m}$. Déterminer $R(t)$ en
    fonction de $R_0$, $t$ et $a$.
    Représenter graphiquement l'allure de la trajectoire de l'une des deux étoiles dans le référentiel
    centré sur le barycentre des deux étoiles.
    \sol{
    On résout en écrivant
    $
    R^3 \d R = -\alpha \d t
    $
    donc $\frac{1}{4}(R^4-R_0^4) = -\alpha t $ et donc $R(t) = R_0(1-\frac{4\alpha t}{R_0^4})^{1/4}$.
    }
    \end{quest}

    Les deux étoiles à neutron sont assimilées à des sphères de diamètre (très faible) $a = \SI{20}{km}$.

    \begin{quest}
    \item Déterminer, en fonction de $R_0$, $a$ et $\alpha$ le temps $t_c$ au bout duquel les deux étoiles entrent en
    contact. Exprimer en fonction de $G$, $M$ et $a$, la vitesse angulaire de rotation
    $\omega_c$ atteinte par le
    système à l'instant $t_c$.
    Application numérique : Calculer $t_c$ et $\omega_c$.
    \sol{
    On a contact lorsque $R(t_c)=a$, soit pour $t_c =\frac{R_0^4}{4\alpha}(1-(a/R_0)^4) = \SI{163}{s}$, et $\omega_c = \sqrt{\frac{2GM}{a^3}} = \SI{6833}{rad.s^{-1}}$.
    }

    \item Justifier que le modèle précédent n'est valable que si la condition $|\dd{R}{t}|\ll R {\omega}$
     est réalisée.
    Cette condition est-elle vérifiée jusqu'à l'instant de contact ?

    \sol{
    Il faut que l'hypothèse que la trajectoire est circulaire soit bien vérifiée.
    Au moment du contact, on a $\dd{R}{t} = \SI{8575}{km.s^{-1}}$ alors que $R \omega_c \simeq \SI{1.4e5}{km.s^{-1}}$. L'hypothèse est donc bien vérifiée, même au contact.
    }

    \end{quest}

\end{exo}
%%%%%%%%%%%%%%%%%%%%%%%%%%%%%%%%%%%

%%%%%%%%%%%%%%%%%%%%%%%%%%%%%%%%%%%
\begin{exo}[PC][1][devoir]{De l'évolution du concept d'atome au cours du \textsc{xx}$^{\rm e}$ siècle}

Ce problème aborde certaines étapes de l'histoire des sciences qui
ont permis, au cours du \textsc{xx}$^{\rm e}$ siècle, de préciser la structure et les
propriétés de l'atome. On s'intéressera
à l'expérience de E. Rutherford, qui conduisit à abandonner le modèle
de J. J. Thomson au profit d'un modèle planétaire.


\subsubsection*{Données numériques}

%Constante de Planck $h=\SI{6,6e-34}{J.s}$

Charge électrique élémentaire $e=\SI{1,6e-19}{C}$

Masse de l'électron $m_{e}=\SI{9,1e-31}{kg}$

Électronvolt $\SI{1}{eV}=\SI{1,6e-19}{J}$

Permittivité du vide $\varepsilon_{0}=\SI{8,9e-12}{A^{2}.s^{4}.kg^{-1}.m^{-3}}$

Célérité de la lumière dans le vide $c=\SI{3,0e8}{m.s^{-1}}$


\subsection{Limite du modèle de J. J. Thomson à travers l'expérience de E. Rutherford}

En 1898, J J Thomson fait l'hypothèse que les atomes sont constitués
d'électrons emprisonnés dans une sorte de gelée de charges positives.
%Ce modèle est appelé modèle du \textquotedbl{}plum pudding\textquotedbl{}, car J J Thomson compare les électrons aux raisins du célèbre dessert anglais.
%%Le physicien Jean Perrin imagine, quant à lui, l'atome à
%l'image du système solaire. Il suppose que les électrons gravitent,
%à des distances immenses, autour d'un «soleil» d'électricité positive,
%sur des orbites pour lesquelles force coulombienne et force d'inertie
%s'équilibrent.

En 1909, Ernest Rutherford, procède à une série d'expériences dans
lesquelles un faisceau de particules alpha (noyaux d'hélium 4: $_{2}^{4}\mbox{He}$),
ayant toutes la même énergie cinétique, est lancé contre une mince
feuille d'or. Il observe que la majorité des particules alpha traversent
la feuille d'or, mais qu'une faible proportion d'entre elles «rebondit»
sur celle-ci. Le but de cette partie est de déterminer se le modèle de Thompson est en accord avec les observations de E. Rutherford.

%Nous nous plaçons d'abord dans le cadre du modèle de J J Thomson,
%supposant une répartition uniforme de la charge positive dans la feuille
%d'or.
%\begin{quest}
%\item  Expliquer qualitativement pourquoi le modèle proposé par J J Thomson
%est incompatible avec les observations de E. Rutherford.
%\end{quest}
%Nous nous plaçons maintenant dans le cadre du modèle de J Perrin,
%supposant l'existence d'un noyau massif de charge positive, et on
%étudie le mouvement de la particule alpha lors de son passage à proximité
%de ce noyau.

Le noyau d'or, de charge positive ponctuelle $Q=Ze$, supposé ponctuel
et immobile dans le référentiel galiléen du laboratoire, se situe
au point $O$, origine d'un repère cartésien orthonormé $(O,\ex,\ey,\ez)$.

Nous considérons qu'à l'instant initial $t=\SI{0}{s}$, la particule
alpha, de masse $m_{\alpha}$ et de charge électrique $q_{\alpha}=2e$,
vient de « l'infini» avec un mouvement rectiligne uniforme caractérisé
par un vecteur vitesse $\Vec{v_{0}}=\Vec{v}(t=0)=v_{0}\ex$.
On désigne par $b$ la distance du point $O$ à la trajectoire de
la particule à l'infini (\Fig{exp_Ruth}). À chaque instant
$t$, on note $d(t)$ la distance entre la particule alpha et le point
$O$. La particule alpha est donc repérée par le vecteur position
$\Vec{OM}(t)=d(t)\er$, avec $(\er,\et,\ez)$
une base cylindrique locale directe.

Au plus proche du point $O$, la particule alpha est au point $S$,
la distance minimale en ce point est notée $d_{m}$. La particule
alpha est non relativiste. L'expérience a été réalisée sous très faible
pression.
%%%%%%%%%%%%%%%%%%%%%
% \begin{figure}[H]
\begin{fig}[label]
  \centering
  \includegraphics[width=0.5\columnwidth]{fig_CCP_PC19_1}
  %\capt{}
  % \protect\caption{Expérience d'Ernest Rutherford.\label{fig:exp_Ruth}}
\end{fig}
%%%%%%%%%%%%

\begin{quest}
\item  On rappelle l'expression de la force de Coulomb qui s'exerce sur la particule alpha,
\begin{equation*}
\Vec{F} =\frac{Qq_\alpha}{4\pi\varepsilon_0d^2}\er=\frac{K}{d^2}\er,
\end{equation*}
avec $K=\frac{Qq_\alpha}{4\pi \varepsilon_0}$.
Donner l'expression de l'énergie potentielle $E_{p}$ qui y est associée,
en considérant que $\lim_{d\rightarrow+\infty}E_{p}(d)=0$, en fonction de $K$, puis de $e$, $Z$, $d$ et $\varepsilon_{0}$.
%Réécrire ces deux expressions
%en fonction de $K=\dfrac{Z\cdot e^{2}}{2\pi\varepsilon_{0}}$ et $d$.\\
Citer les propriétés de cette force qui permettent d'affirmer que
le moment cinétique $\Vec{L_{O}}$ par rapport au point
$O$ et l'énergie mécanique $E_{M}$ de la particule alpha se conservent.

\sol{
Il s'agit de la force électrostatique entre la charge $q_\alpha = 2e$ de la particule alpha et la charge $Q = Ze$ du noyau d'or, donc d'après le cours $$\resm{\vv F=\frac{2Ze^2}{4\pi\varepsilon_0d^2}\er=\frac{K}{d^2}\er}.$$
C'est une force newtonienne dont l'énergie potentielle $E_p$ s'obtient, avec $r=d$, par  $\d E_p = - \delta W(\Vec{F}) = -\Vec{F} \cdot \Vec{\d OM} = -\frac{K \d r}{r^2}$. Donc $$\resm{E_p(d)=\frac{K}{r}=\frac{2Ze^2}{4\pi\varepsilon_0d}}$$ (la constante est nulle pour avoir la limite demandée).
Cette force est à la fois centrale, ce qui implique que le moment cinétique est conservé, et conservative, ce qui implique la conservation de l'énergie mécanique (car c'est la seule force en présence ici).
}

\item  Déterminer, en fonction de $m_{\alpha}$ et $v_{0}$, l'énergie
mécanique $E_{M}$ de la particule alpha.

\sol{
L'énergie mécanique étant constante, elle vaut sa valeur initiale. Or à $t=0$, la particule est à l'infini donc son énergie potentielle est nulle. Il reste donc simplement $\resm{E_M=E_M(t=0)=\frac 1 2 m_\alpha v_0^2}$.
}

\item  Exprimer le moment cinétique $\Vec{L_{O}}$ en fonction
de $b$, $m_{\alpha}$, $v_{0}$ et l'un des vecteurs unitaires du
trièdre direct $(\ex,\ey,\ez)$. Pour cela,
vous pourrez calculer $\Vec{L_{O}}$ en $M_{0}$, position
initiale de la particule alpha telle que $\Vec{OM_{0}}=X\ex+b\,\ey$.
\sol{
Par définition $\vv L_O=\vv{OM_0}\wedge m_\alpha \vv{v_0}=(X\ex+b\ey)\wedge m_\alpha v_0\ex$

\res{\vv L_O=-m_\alpha bv_0\ez}
}
\item  Établir, à un instant $t$ quelconque, l'expression du moment cinétique
$\Vec{L_{O}}$ en fonction de $\dot{\theta}=\dfrac{\mathrm{d}\theta}{\mathrm{d}t}$,
$m_{\alpha}$, $d$ et de l'un des vecteurs unitaires du trièdre direct
$(\ex,\ey,\ez)$.\\
En déduire une relation entre $d$, $b$, $\dot{\theta}$ et $v_{0}$.


\sol{
De même à un instant quelconque : $\vv L_O=\vv{OM}\wedge m_\alpha \vv{v}=d\er\wedge m_\alpha (\dot d\er +d\dot\theta\et)$

\res{\vv L_O=m_\alpha d^2\dot\theta\ez}
On a donc à chaque instant, par conservation de $\vv L_O$ : $\resm{d^2\dot\theta=-bv_0}$.
}

\item  Au sommet $S$ de la trajectoire, le vecteur vitesse $\Vec{v_{s}}$,
de norme $v_{s}$, de la particule alpha est perpendiculaire au rayon
vecteur $\Vec{OS}$, de norme $d_{m}$. Déterminer un polynôme
du second degré en $d_{m}$ et en déduire l'expression de $d_{m}$
en fonction de $K$, $b$, $m_{\alpha}$ et $v_{0}$.
\sol{
Au sommet de la trajectoire, on peut à nouveau calculer $\vv L_O$ facilement car la vitesse est orthogonale au rayon. On en déduit $d_mv_S=bv_0$ par conservation.

Et par conservation de l'énergie mécanique, on a

$$\frac 1 2 m_\alpha v_S^2+\frac K {d_m}=\frac 1 2 m_\alpha v_0^2$$

En remplaçant $v_S$ par la relation précédente et en multipliant tout par $\frac{2d_m^2}{m_\alpha v_0^2}$, on obtient :

\res{d_m^2-\frac{2K}{m_\alpha v_0^2}d_m-b^2=0}

On résout cette équation du second degré :

$$d_m=\frac{K}{m_\alpha v_0^2}\pm\sqrt{\frac{K^2}{m_\alpha^2 v_0^4}+b^2}$$.

Seule la solution avec un "+" est positive donc on garde $\resm{d_m=\frac{K}{m_\alpha v_0^2}+\sqrt{\frac{K^2}{m_\alpha^2 v_0^4}+b^2}}$.
}

\item  Malheureusement, $b$ est inaccessible à la mesure. Par contre,
l'angle de déviation $\varphi$ est facilement mesurable. Il faut
donc trouver la relation qui lie $\varphi$ à $b$. Pour cela, vous
écrirez le principe fondamental de la dynamique (P.F.D.) en fonction
de $K$, $d$, $m_{\alpha}$, $\Vec{v}$ et $\er$. Projeter
le P.F.D. sur l'axe des $x$ en introduisant la composante $v_{x}$
de la vitesse selon l'axe des $x$, et l'angle $\theta$ (\Fig{exp_Ruth}).


Réécrire cette équation en fonction uniquement de $v_{x}$, $\theta$,
$\dot{\theta}$, $K$, $b$, $m_{\alpha}$ et $v_{0}$.


Intégrer cette équation entre $t=0$ et $t\rightarrow\infty$. On
remarquera que $\lim_{t\rightarrow+\infty}\theta(t)\approx\varphi$.


En déduire que la relation qui lie $\varphi$ à $b$ est
$$\tan\left(\dfrac{\varphi}{2}\right)=\dfrac{K}{b\cdot m_{\alpha}\cdot v_{0}^{2}}.$$

\sol{
Le PFD projeté sur $\ex$ donne :

$$m_\alpha \dd{v_x}{t}=\frac{K}{d^2}\cos\theta=-\frac{K}{bv_0}\dot\theta \cos\theta$$

On intègre cette relation entre $t=0$ et $t\to\infty$ :

$$m_\alpha (v_x(\infty)-v_0)=-\frac{K}{bv_0}\left[\sin\theta\right]_0^\infty$$

Or $v_x(\infty)=v_0\cos\varphi$, $\theta(0)=\pi$ et $\theta(\infty)=\varphi$ :

$$m_\alpha v_0(1-\cos\varphi)=-\frac{K}{bv_0}\sin\varphi$$

En utilisant les formules de l'énoncé, on obtient :

$$2m_\alpha v_0\sin^2\frac\varphi 2=\frac{K}{bv_0}\sin\frac\varphi 2\cos\frac\varphi 2$$

Donc en simplifiant on obtient la relation demandée :

\res{\tan \frac\varphi 2=\frac{K}{m_\alpha bv_0^2}}
}

\end{quest}
On rappelle que: $\cos\varphi-1=-2\cdot\sin^{2}\left(\dfrac{\varphi}{2}\right)$
et $\sin\varphi=2\cdot\cos\left(\dfrac{\varphi}{2}\right)\cdot\sin\left(\dfrac{\varphi}{2}\right)$.
\begin{quest}
\item  À partir de quelle valeur de $\varphi$ les particules alpha rebondissent-elles
sur la feuille d'or?
%Expliquer pourquoi le modèle de J Perrin permet
%d'interpréter les observations de E. Rutherford.
\sol{
 Il y aura "rebond" si $\varphi>\frac\pi 2$, pour que les particules repartent en arrière.

Ceci nécessite que $b<\frac{K}{m_\alpha v_0^2}$, donc que $b$ soit "suffisamment petit" : il faut envoyer la particule alpha "suffisamment près" du noyau pour voir l'effet.
}
\end{quest}
Nous nous proposons maintenant d'évaluer une borne supérieure à la
dimension de ce noyau.
\begin{quest}
\item  Montrer que la relation qui lie $d_{m}$ à $\varphi$ est
$$d_{m}=\dfrac{K}{m_{\alpha}v_{0}^{2}}\left(1+\dfrac{1}{\sin\left(\dfrac{\varphi}{2}\right)}\right).$$

\sol{
En reprenant l'expression de $d_m$ :


$$d_m=\frac{K}{m_\alpha v_0^2}\left(1+\sqrt{1+\frac{m_\alpha^2 v_0^4b^2}{K^2}}\right)$$

On reconnait sous la racine l'expression correspondant à $\tan \frac\varphi 2$ :



$$d_m=\frac{K}{m_\alpha v_0^2}\left(1+\sqrt{1+\tan^{-2}\frac\varphi 2}\right)=\frac{K}{m_\alpha v_0^2}\left(1+\sqrt{\frac{\sin^2\frac\varphi 2+\cos^2\frac\varphi 2}{\sin^\frac\varphi 2}}\right)$$

En simplifiant (car $\sin\frac\varphi 2>0$ ici), on obtient bien :

\res{d_m=\frac{K}{m_\alpha v_0^2}\left(1+\frac{1}{\sin\frac\varphi 2}\right)}
}


\item  Pour quelle valeur $\varphi_{m}$ de l'angle $\varphi$, la distance
d'approche est-elle minimale? Déterminer, dans ce cas, l'expression
de $d_{m}$ en fonction de $K$, $m_{\alpha}$ et $v_{0}$.

\sol{
$d_m$ est minimale lorsque $\sin\frac\varphi 2$ est maximal donc pour $\varphi=\resm{\varphi_m= \pi}$. On a alors $\resm{d_{m, min}=\frac{2K}{m_\alpha v_0^2}}$.
}


\item  Que vaut $b$ pour $\varphi=\varphi_{m}$? Représenter l'allure
de la trajectoire de la particule alpha pour cet angle et faire figurer
$d_{m}$ sur votre schéma. Justifier que $d_{m}$ constitue une borne
supérieure du rayon du noyau.



Sachant que l'énergie typique d'une particule alpha est de $\SI{5}{MeV}$
et que le numéro atomique de l'or est $Z=79$, déterminer numériquement
la valeur de $d_{m}$.

\sol{
Avec la relation de la \item  , on vérifie que cette situation correspond à $\resm{b\to0}$. La trajectoire est alors rectiligne, droit en direction du noyau, avec un demi-tour à une distance $d_{m, min}$ de $O$. Si la particule peut arriver jusqu'en $d_{m, min}$ avant de faire demi-tour, c'est que le noyau est un peu plus loin, donc plus petit que $d_{m, min}$ .


 $d_{m, min}=\frac{2K}{m_\alpha v_0^2}=\frac{K}{E_c}=\frac{Ze^2}{2\pi\varepsilon_0 E_c}\approx 5.10^{-14}$ m.
}

\item  Justifier que, pour effectuer des expériences de physique nucléaire,
il faut disposer de particules de haute énergie.

\sol{
Pour sonder l'intérieur du noyau, il faut diminuer $d_{m, min}$ donc augmenter l'énergie cinétique des particules incidentes.
}
\end{quest}

\subsection{Limite du modèle planétaire}

Le modèle de J. J. Thomson est écarté et l'on considère que les électrons
évoluent, avec un mouvement circulaire uniforme, autour d'un noyau
massif de charge électrique positive. Néanmoins, ce modèle est en
contradiction avec une loi classique de l'électromagnétisme: toute
particule chargée et accélérée émet de l'énergie électromagnétique.

Pour mettre en évidence les conséquences de cette loi classique de
l'électromagnétisme, nous allons étudier le mouvement de l'électron
de l'atome d'hydrogène, de masse $m_{e}$ et de charge électrique
$q_{e}=-e$, qui tourne autour de son noyau, un proton de masse $m_{p}$
et de charge électrique $q_{p}=+e$, sur une orbite circulaire de
rayon $r$ (\Fig{Modele-planetaire}). Le noyau est considéré,
dans le référentiel galiléen du laboratoire, fixe, ponctuel et placé
en son centre $C$. Le centre de la trajectoire circulaire de l'électron
est donc $C$.

% \begin{figure}[H]
\begin{fig}[label]
  \centering
  \includegraphics[width=0.4\columnwidth]{fig_CCP_PC19_2}
  % \protect\caption{Modèle planétaire de l'atome d'hydrogène.\label{fig:Modele-planetaire}}
\end{fig}



Pour étudier le mouvement circulaire de l'électron, nous allons utiliser
le repère polaire pour lequel, en un point $M$ de la trajectoire
décrite par l'électron, on associe deux vecteurs unitaires $\er$
et $\et$ (\Fig{Modele-planetaire}). $\et$
est le vecteur tangent à la trajectoire au point $M$ et dirigé dans
le sens du mouvement. La position de 1'électron est repérée par le
vecteur position: $\Vec{CM}=r\er$ et l'angle
$\theta=(\Vec{Cx},\Vec{CM})$.
\begin{quest}
\item  Déterminer l'expression du vecteur vitesse $\Vec{v}$ de l'électron
en fonction de $e$, $m_{e}$, $\varepsilon_{0}$, $r$ et d'un vecteur
unitaire.

\sol{
On applique le PFD à l'électron, soumis uniquement à la force électrostatique et on le projette dans la base cylindrique (avec $r=cste$) :

$$m_e\dd{v}{t}=-\frac{e^2}{4\pi\varepsilon_0 r^2}\er$$

$$\left\{ \begin{array}{l}
-m_er\dot\theta^2=-\frac{e^2}{4\pi\varepsilon_0 r^2}\\
m_er\ddot\theta = 0
\end{array}\right.$$

On en déduit que la norme de la vitesse est constante et que finalement $\resm{\vv v=r\dot\theta\et=\frac{e}{\sqrt{4\pi\varepsilon_0 m_e r}}\et}$.


}
\item  Exprimer l'énergie mécanique $E_{M}(r)$ de l'électron sous la forme
$E_{M}(r)=A\cdot f(r)$ où $A$ est une constante négative dont vous
préciserez l'expression en fonction de $e$, $\varepsilon_{0}$ et
$f(r)$ une fonction qui ne dépend que de $r$ que vous déterminerez
également.

\sol{
L'énergie mécanique s'écrit : $E_M=\frac{1}{2} m_e v^2 - \frac{e^2}{4\pi\varepsilon_0 r} = -\frac{e^2}{8\pi\varepsilon_0 r}$ en remplaçant par l'expression de $v$. On a donc $A=-\frac{e^2}{8\pi\varepsilon_0}$ et $f(r)=\frac{1}{r}$.
}

\item  Une loi classique de l'électromagnétisme indique que toute particule
chargée et accélérée émet de l'énergie électromagnétique. Aussi, d'après
cette théorie, l'électron devrait émettre un rayonnement électromagnétique
de puissance moyenne:
\[
P(r)=\dfrac{\omega^{4}\cdot e^{2}\cdot r^{2}}{12\cdot\pi\cdot\varepsilon_{0}\cdot c^{3}}
\]
où $\omega$ est la vitesse angulaire de l'électron et $c$ la vitesse
de la lumière dans le vide.


Cette puissance peut être mise sous la forme $P(r)=P_{0}\dfrac{1}{r^{4}}$,
où $P_{0}$ est une constante.


Déterminer l'expression de $P_{0}$ et son unité.


Justifier que le rayon de la trajectoire de l'électron diminue au
cours du temps.

\sol{
$\omega=\dot\theta=\frac{e}{\sqrt{4\pi\varepsilon_0 m_e r^3}}$ donc en remplaçant dans l'expression de $P(r)$, on obtient :

\res{P(r)=\frac{P_0}{r^4}\textbf{ avec }P_0=\frac{e^6}{192\pi^3c^3\varepsilon_0^3m_e^2}}

Si $P(r)$ est une puissance (\textcolor{red}{fallait-il le redémontrer ?}), $P_0$ s'exprime en W.m$^{4}$.

$P(r)$ représente l'énergie perdue par rayonnement par unité de temps, donc l'énergie mécanique de l'électron diminue. Comme $A<0$ et $f(r)$ est décroissante, le rayon diminue aussi.
}

\item  Montrer qu'il existe une relation différentielle de la forme: $r^{2}\cdot\dfrac{\mathrm{d}r}{\mathrm{d}t}=\dfrac{P_{0}}{A}$.

\sol{
On applique le théorème de la puissance mécanique (ie un bilan d'énergie mécanique) à l'électron, sachant que $P(r)$ est la puissance perdue. On a donc :

$$\dd{E_M}{t}=-P(r)\Rightarrow -\frac A {r^2}\dd{r}{t}=-\frac{P_0}{r^4}$$

\res{r^2\dd{r}{t}=\frac{P_0}{A}}

Comme $A<0$, le rayon diminue bien.
}

\item  À $t=0$, on suppose que l'électron se trouve sur une orbite de
rayon $R$. Donner l'expression, en fonction de $P_{0}$, $R$ et
$A$, du temps $t_{f}$ mis par l'électron pour atteindre le noyau.


On donne $R=\SI{1,0e-10}{m}$, calculer $t_{f}$. Commenter le résultat
obtenu.

\sol{
On sépare les variables avant d'intégrer :

$$\int_R^0 r^2\d r=\int_0^{t_f}\frac{P_0}{A}\d t\Rightarrow \resm{t_f=-\frac{AR^3}{3P_0}}$$

 $\frac{-A}{P_0}=\frac{24\pi^2c^3\varepsilon_0^2m_e^2}{e^4}\approx 6,4.10^{20}$ s.m$^{-3}$ donc $t_f\approx 2,1.10^{-10}$ s. L'atome ne serait stable que pendant un temps infiniment court, c'est impossible.
}

\end{quest}

\end{exo}
%%%%%%%%%%%%%%%%%%%%%%%%%%%%%%%%%%%



%%%%%%%%%%%%%%%%%%%%%%%%%%%%%%%%%%%%%%%%%%%
\begin{exo}[MPSI][1][colle]{Tir au basketball}
  \Notions{\item Dynamique du point \task PFD \task Chute libre \task  Trajectoire \task Coordonnées cartésiennes}
 Un joueur de basketball est placé à une distance $D$ du panier qui est à une hauteur~$H$. Il tire le ballon d'une hauteur $h$ en
 lui imposant une vitesse $\Vec{v_{0}}$ faisant un angle~$\alpha$ avec l'horizontale. On néglige les frottements. On note $\Vec{e_{x}}$ l'horizontale, et
 $\Vec{e_{z}}$ la verticale pointant vers le haut.
 \begin{quest}
  \item Appliquer le principe fondamental de la dynamique au ballon, puis déterminer les équations horaires $x(t)$ et $z(t)$.

  \sol{

  \begin{liste}
    \item \udl{Système} : $\{$ballon$\}$ assimilé à une point matériel $\rm M$.
    \item \udl{Référentiel} : référentiel du sol $\Rcal$, supposé galiléen.
    \item \udl{Actions mécaniques} : 
    % \begin{itemize}
      % \item 
      poids : $\Vec{P} = m\Vec{g}$, avec $\Vec{g} = -g\Vec{e_z}$.
    % \end{itemize}

    \item \udl{Principe fondamental de la dynamique (PFD)} :
    $$
    m \aR(M) = \Vec{P} = m\Vec{g} = - mg \ez.
    $$
  \end{liste}
  Par projection sur la base cartésienne $(\ex,\ez)$, on obtient :
    \begin{equation*}
      \begin{cases}
        \ddot{x} = 0 ,\\
   \displaystyle \ddot{z} = -g.
  \end{cases}
 \end{equation*}

  On intègre une première fois en prenant en compte la condition initiale sur la vitesse : $\vR(t=0) = v_0 (\cos \alpha \, \ex + \sin \alpha \,\ez)$,
$$
\begin{cases}
\dot{x} = v_{0}\,\cos(\alpha), \\
\displaystyle \dot{z} = -gt + v_{0}\,\sin(\alpha).
\end{cases}
$$
Puis on intègre une seconde fois en prenant en compte la condition initiale sur la position : $x(0) = 0$ et $z(0) = h$,

      % \begin{equation*}
      \eqbox{
        \begin{cases}
          x = v_{0}\,\cos(\alpha)\,t, \\
          \displaystyle z = -\frac{g}{2}t^{2} + v_{0}\,\sin(\alpha)\,t + h.
        \end{cases}
        }
      %  \end{equation*}
  }

  \item En déduire l'équation de la trajectoire du ballon $z(x)$.

  \sol{
  Par définition, la \udl{trajectoire} du ballon est l'~\fbox{équation $z(x)$} reliant les coordonnées $x$ et $z$ du ballon.
  On peut proposer deux méthodes pour trouver la trajectoire.
  \begin{enumerate}
    \item Méthode \guill{classique} : on inverse la relation $x(t)$ en $t(x)$, en vérifiant que la fonction $x(t)$ est bien bijective, puis en remplace $t$ dans l'expression de $z(t)$. D'après les lois horaires trouvées précédemment, on a 
    $$
    t(x) = \frac{x}{v_{0}\,\cos(\alpha)}.
    $$
    puis
    $$
    z = -\frac{g}{2}t^{2} + v_{0}\,\sin(\alpha)\,t + h = -\frac{g}{2}\left(\frac{x}{v_{0}\,\cos(\alpha)}\right)^{2} + v_{0}\,\sin(\alpha)\left(\frac{x}{v_{0}\,\cos(\alpha)}\right) + h,
    $$
    soit 
    \eqbox{
    z(x) = -\frac{g}{2v_{0}^{2}\cos^{2}(\alpha)}\,x^{2} + \tan(\alpha)\,x + h.
    }
    \item Méthode \guill{originale} : on détermine la dérivée $z'(x)$ à partir de la formule de dérivation en chaîne, puis on intègre pour trouver $z(x)$.
    $$
    \dd{z}{x} = \dd{z}{t} \dd{t}{x} = \frac{\dot{z}}{\dot{x}} = \frac{-gt + v_{0}\,\sin(\alpha)}{v_{0}\,\cos(\alpha)} = -\frac{g}{v_{0}\,\cos(\alpha)}\,t + \tan(\alpha) = -\frac{g}{v_{0}^{2}\cos^{2}(\alpha)}\,x + \tan(\alpha).
    $$
    En intégrant puis en utilisant la condition initiale sur $z(x)$, on trouve bien l'équation de trajectoire précédente. Cette méthode ne nécessite pas de connaître $z(t)$, mais il faut aussi connaître $t(x)$.
  \end{enumerate}
  % Trajectoire $\displaystyle z(x) = -\frac{g}{2v_{0}^{2}\cos^{2}(\alpha)}\,x^{2} + \tan(\alpha)\,x + h$.}
  }

  \item Quelle équation doit vérifier $\tan(\alpha)$ pour que le panier soit marqué ? En déduire une condition sur la norme de la vitesse initiale $v_{0}$ pour que cette équation admette une solution. Cette condition est-elle vérifiée pour $v_0 = \SI{1}{m.s^{-1}}$ ? $v_0 = \SI{3}{m.s^{-1}}$ ? On donne $D=\SI{5.00}{m}$, $H=\SI{3.05}{m}$, $h=\SI{2.00}{m}$ ?

  \sol{

  Pour marquer le panier, la trajectoire doit vérifier l'équation 
  \eqbox{
    $z(D)=H$.
  }
  Cette condition se récrit comme une équation algébrique du second ordre sur le paramètre $X = \tan \alpha$ :

  $$
  -\frac{g}{2v_{0}^{2}\cos^{2}(\alpha)}\,D^{2} + \tan(\alpha)\,D + h = H,
  $$
  soit  avec $\cos(\alpha)^{-2} = 1+ \tan(\alpha)^{2}$ :
      % \begin{equation*}
      \eqbox{
        X^2 - \frac{2v_{0}^{2}}{gD} X - \frac{2(h-H)v_{0}^{2}}{gD^{2}} +1 = 0.
        }
      % \end{equation*}
   On impose que le discriminant de l'équation précédente soit positif, pour qu'il existe une solution $X$ réelle,
 \begin{equation}
  v_{0}^{4}-2g(H-h)\,v_{0}^{2}-g^{2}D^{2} \geq 0.
 \end{equation}
 On trouve là-aussi une condition sur la grandeur $v_0^2$ sous forme d'une équation du second degré. La racine positive du polynôme précédent vaut 
 $$
 v_{\rm min}^2 = g\left(H-h+\sqrt{(H-h)^2+D^2}\right).
 $$
 La vitesse initiale doit être supérieure à $v_{\rm min}$ pour que le panier puisse être marqué, à condition également que le ballon soit envoyé avec le bon angle $\alpha$. On a donc
 \eqbox{
 v_0 \geq \sqrt{g\left(H-h+\sqrt{(H-h)^2+D^2}\right)} \simeq \SI{4.0}{m.s^{-1}}.
 }
  }

 \end{quest}

%  $\begin{array}{lll}
% \textsc{Données :}
% & \text{distance du panier} & : D=\SI{5.00}{m} \\
% & \text{hauteur du panier} & :H=\SI{3.05}{m}\\
% & \text{hauteur du tir} & : h=\SI{2.00}{m}
% \end{array}
% $

\end{exo}
%%%%%%%%%%%%%%%%%%%%%%%%%%%%%%%%%%%%%%%%

%%%%%%%%%%%%%%%%%%%%%%%%%%%%%%%%%%%%%%%%%%%%%%%%%%%%%%%%%%%%%%%%
\begin{exo}[PCSI][1][colle]{Balistique d'un projectile dans l'air}
On considère la situation d'un tir d'obus assimilé à un point matériel de masse $m$ dans le champ de pesanteur $\vv{g}$ dirigé selon $-\ez$. Dans un premier temps on néglige l'influence de la résistance de l'air, c'est-à-dire qu'on néglige le freinage aérodynamique. Les conditions initiales qui définissent le problème du tir sont $\vv{r_0}=0$, $\vv{v_0}=v_0\left(\cos\alpha\,\ex+\sin\alpha\,\ez\right)$.
\begin{quest}
\item Déterminer les équations horaires du mouvement, et en déduire la nature de la trajectoire puis la portée du tir.
\item Déterminer la valeur de l'angle $\alpha_0$ qui rend la portée maximale pour un canon donné, puis la valeur de ce maximum. Montrer que pour une valeur inférieure à la portée maximale, un point du plan $z=0$ peut être atteint par le canon de deux façons dont on précisera comment les inclinaisons $\alpha_1$ et $\alpha_2$ sont disposées par rapport à $\alpha_0$.
\item On reprend maintenant le problème en modélisant la résistance de l'air par une force $-m\lambda\vv{v}$. Déterminer à nouveau les équations du mouvement et préciser les différences avec le cas précédent pour $t\rightarrow\infty$.
\end{quest}
\end{exo}
%%%%%%%%%%%%%%%%%%%%%%%%%%%%%%%%%%%%%%%%%%%%%%%%%%%%%%%%%%%%%%%%


%%%%%%%%%%%%%%%%%%%%%%%%%%%%%%%%%%%
\begin{exo}[PC][1][colle]{Vitesse d'un satellite}
La vitesse instantanée $v$ d'un satellite en orbite autour d'un astre à la distance $r$ est donnée par la formule suivante :
\eq{
v(r) = \frac{\alpha}{r^{\beta}}.
}

On dispose de plus de certaines données :
\begin{center}
	\begin{tabular}{c||c c c c c c }
	    $r$ (km) & 500  & 1000 & 2000 & 5000 & 10000 & 42000 \\
	\hline
	$v$ ($\SI{}{km.s^{-1}}$) & 28 & 20 & 14 & 8.9 & 6.3 & 3.1 \\
	\end{tabular}
\end{center}

\begin{quest}
\QR{
Déduire le signe de $\beta$ à partir de l'étude des variations de $v$ en fonction de $r$.
}{}

\QR{
Linéariser le modèle proposé (à l'aide de la fonction logarithme).
}{}

\QR{
Obtenir $\alpha$ et $\beta$ à l'aide de votre calculette.
}{}

\QR{
Que vaut alors la dimension de $\alpha$ ? Expliquer cette loi.
%Quelle est son expression d'après les équations de la mécanique ?
}{}
\end{quest}

\end{exo}
%%%%%%%%%%%%%%%%%%%%%%%%%%%%%%%%%%%

%%%%%%%%%%%%%%%%%%%%%%%%%%%%%%%%%%%
\begin{exo}[PC][1][colle]{Pendule simple}
\Notions{\item Dynamique du point \task Pendule simple \task TMC \task Coordonnées polaires}
	\noindent	 Une masse ponctuelle $m$ est suspendue à l'extrémité P d'un fil inextensible de masse négligeable et de longueur $L=OP$. Il peut effectuer des mouvements de rotation dans le plan vertical ($Oxy$), autour de l'axe horizontal ($Oz$). La position du pendule est repérée par l'angle $\theta$ que fait le fil avec la verticale.

\begin{fig}[]
  \centering
  \includegraphics[width=0.3\textwidth]{schema_pendule_simple2.png}
\end{fig}

		 L'étude est menée dans le référentiel terrestre considéré comme galiléen. Les frottements sont négligés. L'ensemble ainsi décrit se trouve dans le champ de pesanteur terrestre d’accélération $\Vec{g}=g\ex$.

		%  On observe que la période des oscillations d’un pendule simple lorsque son amplitude est petite est :
		% \eq{
		% T=2\pi \sqrt{\frac{L}{g}}.
		% }


\begin{quest}

% \QR{
% 	 Vérifier que cette formule est homogène.
% }{
% 	 dim(L)=L, dim(g)=L/T$^2$, dim(T)=T
%
% }

\QR{
	 Exprimer les forces qui s’appliquent sur la masse $m$ dans la base polaire $(\er,\et)$.
}{
	 poids et tension du fil
}

\item On rappelle l’expression de la vitesse $\Vec{v}=L \dot{\theta} \et$ et de l’accélération $\Vec{a}=L\ddot{\theta}\et - L \dot{\theta}^2 \er$ en coordonnées polaires. Établir l’équation différentielle vérifiée par $\theta(t)$.
% \end{quest}
% 		L’équation différentielle du mouvement s’écrit :
% 		\eq{
% 		\frac{d^2 \theta}{dt^2}+\frac{g}{L}\sin (\theta) = 0
% 		}
%
% \begin{quest}
\QR{
Simplifier cette équation en supposant $\theta \ll 1$. Quelle est le nom de
l’équation ainsi obtenue ? Exprimer la période $T$ des petites oscillations et vérifier l’homogénéité de cette formule. \AN~pour $L=\SI{20}{cm}$.
}{
	 Le système constitue un oscillateur harmonique. $T=2\pi \sqrt{\frac{L}{g}}$.
   }

% \end{quest}
    %Lorsque $\theta_0$ est petit, $\sin (\theta) \approx \theta$.
% \begin{quest}
\QR{
\`A l'instant initial, la vitesse est nulle et $\theta(t=0)=\theta_0$.
	 Exprimer la solution de l’équation différentielle compatible avec les conditions initiales.
}{

	\eq{
	\frac{d^2 \theta}{dt^2}+\omega_0^2 \theta = 0
	\qquad \Rightarrow \qquad \omega_0=\sqrt{\frac{g}{L}} \qquad \& \qquad \theta(t=0)=\theta_0 \quad \& \quad \frac{d \theta}{dt}(t=0)=0.
	}

	\eq{
	\boxed{\theta(t)=\theta_0 \cos (\omega_0 t)}.
	}
}
\end{quest}
\end{exo}
%%%%%%%%%%%%%%%%%%%%%%%%%%%%%%%%%%%


%%%%%%%%%%%%%%%%%%%%%%%%%%%%%%%%%%%
\begin{exo}[PC][1][colle]{Moment cinétique et pendule simple amorti}
Un objet ponctuel A de masse m est suspendu à l'extrémité P d'un fil inextensible de masse négligeable et de longueur $L=OP$. Il peut effectuer des mouvements de rotation dans le plan vertical ($Oxy$), autour de l'axe horizontal ($Oz$). La position de l'objet A est repérée par l'angle $\theta$ que fait le fil avec la verticale.

\begin{fig}[]
  \centering
  \includegraphics[width=0.3\textwidth]{schema_pendule_simple2.png}
\end{fig}

		 L'étude est menée dans le référentiel terrestre considéré comme galiléen.L'ensemble ainsi décrit se trouve dans le champ de pesanteur terrestre d’accélération $\Vec{g}=g\ex$.

		%  On observe que la période des oscillations d’un pendule simple lorsque son amplitude est petite est :
		% \eq{
		% T=2\pi \sqrt{\frac{L}{g}}.
		% }


\begin{quest}

% \QR{
% 	 Vérifier que cette formule est homogène.
% }{
% 	 dim(L)=L, dim(g)=L/T$^2$, dim(T)=T
%
% }

\QR{
	 Exprimer les moments des forces en $O$ qui s’appliquent sur la masse $m$ en négligeant les frottements.
}{
	 poids et tension du fil
}

\item On rappelle l’expression de la vitesse $\Vec{v}=L \dot{\theta} \et$ en coordonnées polaires. Calculer le moment cinétique $\Vec{L}_O$ et appliquer le théorème du moment cinétique.
% \end{quest}
% 		L’équation différentielle du mouvement s’écrit :
% 		\eq{
% 		\frac{d^2 \theta}{dt^2}+\frac{g}{L}\sin (\theta) = 0
% 		}
%
% \begin{quest}
\QR{
En déduire l'équation différentielle vérifiée par $\theta(t)$ en supposant $\theta \ll 1$. Quelle est le nom de
l’équation ainsi obtenue ? Exprimer la période $T$ des petites oscillations. \AN~pour $L=\SI{20}{cm}$.
}{
	 Le système constitue un oscillateur harmonique. $T=2\pi \sqrt{\frac{L}{g}}$.
   }

% \end{quest}
    %Lorsque $\theta_0$ est petit, $\sin (\theta) \approx \theta$.
% \begin{quest}
\QR{
L'air exerce une force de frottement fluide $\Vec{F} = - \lambda \Vec{v}$. Appliquer le théorème du moment cinétique et déterminer l'équation différentielle vérifiée par $\theta(t).$
}{

	\eq{
	\frac{d^2 \theta}{dt^2}+\omega_0^2 \theta = 0
	\qquad \Rightarrow \qquad \omega_0=\sqrt{\frac{g}{L}} \qquad \& \qquad \theta(t=0)=\theta_0 \quad \& \quad \frac{d \theta}{dt}(t=0)=0.
	}

	\eq{
	\boxed{\theta(t)=\theta_0 \cos (\omega_0 t)}.
	}
}
\end{quest}
\end{exo}
%%%%%%%%%%%%%%%%%%%%%%%%%%%%%%%%%%%


%%%%%%%%%%%%%%%%%%%%%%%%%%%%%%%%%%%
\begin{exo}[PC][1][colle]{Approche énergétique du pendule simple}

	\noindent	 Un objet ponctuel A de masse m est suspendu à l'extrémité P d'un fil inextensible de masse négligeable et de longueur $L=OP$. Il peut effectuer des mouvements de rotation dans le plan vertical ($Oxy$), autour de l'axe horizontal ($Oz$). La position de l'objet A est repérée par l'angle $\theta$ que fait le fil avec la verticale.

\begin{fig}[]
  \centering
  \includegraphics[width=0.3\textwidth]{schema_pendule_simple2.png}
\end{fig}

		 L'étude est menée dans le référentiel terrestre considéré comme galiléen. Les frottements sont négligés. L'ensemble ainsi décrit se trouve dans le champ de pesanteur terrestre d’accélération $\Vec{g}=g\ex$.

		%  On observe que la période des oscillations d’un pendule simple lorsque son amplitude est petite est :
		% \eq{
		% T=2\pi \sqrt{\frac{L}{g}}.
		% }


\begin{quest}

% \QR{
% 	 Vérifier que cette formule est homogène.
% }{
% 	 dim(L)=L, dim(g)=L/T$^2$, dim(T)=T
%
% }

\QR{
	  Exprimer les énergies cinétique et potentielle en fonction de $L$, $m$, $\dot{\theta}$, $\theta$ et $g$.
}{
}

\item Que peut-on dire de l'énergie mécanique du système ? En déduire l'équation différentielle vérifiée par l'angle $\theta(t)$.

\QR{
Simplifier cette équation en supposant $\theta \ll 1$. Quelle est le nom de
l’équation ainsi obtenue ? Exprimer la période $T$ des petites oscillations et vérifier l’homogénéité de cette formule. \AN~pour $L=\SI{20}{cm}$.
}{
	 Le système constitue un oscillateur harmonique. $T=2\pi \sqrt{\frac{L}{g}}$.
   }

% \end{quest}
    %Lorsque $\theta_0$ est petit, $\sin (\theta) \approx \theta$.
% \begin{quest}
\QR{
\`A l'instant initial, la vitesse est nulle et $\theta(t=0)=\theta_0$.
	 Exprimer la solution de l’équation différentielle compatible avec les conditions initiales.
}{

	\eq{
	\frac{d^2 \theta}{dt^2}+\omega_0^2 \theta = 0
	\qquad \Rightarrow \qquad \omega_0=\sqrt{\frac{g}{L}} \qquad \& \qquad \theta(t=0)=\theta_0 \quad \& \quad \frac{d \theta}{dt}(t=0)=0.
	}

	\eq{
	\boxed{\theta(t)=\theta_0 \cos (\omega_0 t)}.
	}
}
\end{quest}
\end{exo}
%%%%%%%%%%%%%%%%%%%%%%%%%%%%%%%%%%%

%%%%%%%%%%%%%%%%%%%%%%%%%%%%%%%%%%%%%%%%%%%%%
\begin{exo}[PCSI][1][colle]{Modèle de Drude}
On considère un conducteur électrique assimilé à un cylindre d'axe $(Oz)$ comportant des charges mobiles: les électrons de conduction. À partir de l'instant $t=0$, ces électrons sont mis en mouvement par un champ électrostatique $\vv{E}=E\,\vv{e_z}$. Il sont aussi soumis, à cause des chocs avec le réseau sous-jacent (les cations du conducteur), à une force de frottement de la forme $\vv{f}=-\frac{m}{\tau}\vv{v}$ où $m$ est la masse de l'électron et $\tau$ la durée typique entre deux collisions successives d'un électron avec le réseau.
\begin{quest}
\item Déterminer la vitesse $\vv{v}(t)$ et en déduire l'expression de la vitesse limite $\vv{v_{\rm{lim}}}$.
\item Le nombre d'électrons mobiles par unité de volume est noté $n$. Lorsque le régime permanent est établi, montrer que le vecteur densité volumique de courant $\vv{j}$ se met sous la forme $\vv{j}=\sigma\vv{E}$ où $\sigma$ est appelée la conductivité électrique du métal. Cette expression est appelée la loi d'Ohm locale, vous la reverrez en deuxième année !
\end{quest}
\end{exo}
%%%%%%%%%%%%%%%%%%%%%%%%%%%%%%%%%%%%%%%%%%%%%

%%%%%%%%%%%%%%%%%%%%%%%%%%%%%%%%%%%%%%%%%%%%%
\begin{exo}[PCSI][1][colle]{Particule dans des champs E et B parallèles}
On étudie dans le référentiel terrestre supposé galiléen le mouvement d'une particule chargée de charge $q$ et de masse $m$, initialement de vitesse $\vv{v_0}=v_{0}\,\vv{e_y}$ et à l'origine du repère, dans un champ électromagnétique uniforme et stationnaire $\vv{E}=E\,\vv{e_z}$ et $\vv{B}=B\,\vv{e_z}$.
\begin{quest}
\item Déterminer les équations différentielles régissant le mouvement de la particule. En déduire l'équation horaire $z(t)$ de la trajectoire et le mouvement de la particule en projection selon l'axe $(Oz)$.
\item Déterminer les équations horaires $x(t)$ et $y(t)$ et en déduire la nature du mouvement de la particule en projection dans le plan $(Oxy)$.
\item Conclure sur la nature complète de la trajectoire de la particule dans le référentiel terrestre.
\end{quest}
\end{exo}
%%%%%%%%%%%%%%%%%%%%%%%%%%%%%%%%%%%%%%%%%%%%%

%%%%%%%%%%%%%%%%%%%%%%%%%%%%%%%%%%%%%%%%%%%%%
\begin{exo}[PCSI]{Le toboggan}
Un enfant, assimilé à un point matériel de masse $m$, glisse sur un toboggan en forme de quart de cercle de rayon $r$. Initialement au point $A$, l'enfant se laisse glisser sans vitesse initiale et atteint le point $B$ avec une vitesse $v$. On néglige les frottements.
%\begin{figure}[h!]
% \centering
\begin{center}
 \includegraphics[width=0.5\textwidth]{toboggan.png}
\end{center}
% \end{figure}
\begin{quest}
\item Avec le théorème du moment cinétique, établir l'équation différentielle vérifiée par $\theta(t)$.
\item En multipliant la relation précédente par $\dot{\theta}(t)$ et en intégrant, déterminer $\dot{\theta}(t)$ puis en déduire l'expression de la vitesse au point $B$. Faire l'application numérique pour $r=\SI{2,5}{m}$, $\theta_{0}=\SI{10}{\degree}$ et commenter.
\item Retrouver le résultat précédent par une méthode énergétique.
\end{quest}
\end{exo}
%%%%%%%%%%%%%%%%%%%%%%%%%%%%%%%%%%%%%%%%%%%%%

%%%%%%%%%%%%%%%%%%%%%%%%%%%%%%%%%%%%%%%%%%%%%
\begin{exo}[PCSI]{Modèle classique d'un trou noir}
Cet exercice propose, dans un modèle non relativiste et non quantique, de chercher à déterminer l'ordre de grandeur de la taille d'un trou noir. On considère un point matériel $M$ de masse $m$ situé à proximité d'un astre sphérique de masse $m_0$, de rayon $R$ et de centre $O$. Cet astre est supposé suffisamment massif pour que le point matériel ne soit soumis qu'à la force gravitationnelle due à l'astre. On étudie de mouvement de $M$ dans le référentiel lié à l'astre, supposé galiléen.
\begin{quest}
\item Exprimer la force gravitationnelle ressentie par $M$ ainsi que l'énergie potentielle dont elle dérive en la supposant nulle à l'infini. Exprimer l'énergie mécanique de $M$. Est-elle conservée ?
\item Montrer que le mouvement de $M$ est plan. En se plaçant en coordonnées polaires, justifier que $C=r^2\dot{\theta}$ est une constante du mouvement.
\item Montrer que l'énergie mécanique se met sous la forme $E_{\rm{m}}=\frac{1}{2}m\dot{r}^2+E_{\rm{p,eff}}(r)$ en précisant l'expression de l'énergie potentielle effective, puis tracer son allure en fonction de $r$.
\item Déterminer graphiquement les valeurs de $E_{\rm{m}}$ pour lesquelles le point $M$ peut échapper à l'attraction de l'astre. Comment nomme-t-on ces états ? En déduire la vitesse de libération $v_{\rm{lib}}$ à la surface de l'astre.
\item En supposant qu'un trou noir est un astre dont la vitesse de libération est supérieure à $c=\SI{3e8}{m.s^{-1}}$, calculer le rayon maximal que l'astre doit avoir pour être un trou noir, appelé rayon de Schwarzschild $R_{\rm{S}}$.
\end{quest}
\end{exo}
%%%%%%%%%%%%%%%%%%%%%%%%%%%%%%%%%%%%%%%%%%%%%

%%%%%%%%%%%%%%%%%%%%%%%%%%%%%%%%%%%%%%%%%%%%%
\begin{exo}[PCSI]{Bille dans une cuvette parabolique}
On place une bille sphérique, assimilée à un point matériel de masse $m$, dans une cuvette parabolique (par exemple un saladier) dont l'équation est de la forme $z=\frac{r^2}{a}$ en coordonnées cylindriques. On notera $O$ le point situé au fond du saladier. On suppose que la bille se déplace sans frottement dans la saladier et on étudie le mouvement dans le référentiel terrestre supposé galiléen.
\begin{quest}
\item Exprimer dans la base cylindrique le moment cinétique du point $M$ par rapport à $O$, noté $\vv{L}_O$. Montrer que la composante de $\vv{L}_O$ selon la verticale ascendante $\vv{e_z}$ se conserve au cours du temps. En notant $L$ cette constante, l'exprimer en fonction de $m,r$ et $\dot{\theta}$.
\item Exprimer l'énergie cinétique $E_{\rm{c}}$ du point $M$ dans le référentiel terrestre. Justifier ensuite l'existence d'une énergie potentielle pour $M$, que l'on exprimera en fonction de $r$, en supposant que $E_{\rm{p}}(0)=0$. Que peut-on dire de l'énergie mécanique de $M$ ?
\item Déduire de ce qui précède que $E_{\rm{m}}=\frac{1}{2}m\dot{r}^2G(r)+E_{\rm{p,eff}}(r)$ en précisant l'expression de $G(r)$ et $E_{\rm{p,eff}}$. Représenter le graphe, et discuter de la nature du mouvement de $M$.
\end{quest}
\end{exo}
%%%%%%%%%%%%%%%%%%%%%%%%%%%%%%%%%%%%%%%%%%%%%

%%%%%%%%%%%%%%%%%%%%%%%%%%%%%%%%%%%%%%%%%%%%%
\begin{exo}[PCSI]{Descente à ski}
Un skieur de masse $m=\SI{65}{kg}$ et assimilé à un parallélépipède descend une piste de ski de longueur $l=\SI{50}{m}$ assimilée à un plan incliné de $\alpha=\SI{25}{\degree}$ par rapport à l'horizontale. On notera $\mu=\num{0,1}$ le coefficient de frottement de Coulomb entre les skis et la piste lorsqu'il y a glissement.
\begin{quest}
\item Faire un bilan des forces, donner leurs expressions et relier deux d'entre elles par la loi de Coulomb que l'on rappellera.
\item Exprimer puis calculer numériquement le travail de toutes les forces lors de la descente.
\item En supposant que le skieur part sans vitesse initiale et qu'il descend la piste \guill{tout-schuss} (c'est-à-dire sans prendre aucun virage), quelle sera sa vitesse à l'arrivée ? Convertir la valeur en \SI{}{km.h^{-1}} et commenter.
\end{quest}
\end{exo}
%%%%%%%%%%%%%%%%%%%%%%%%%%%%%%%%%%%%%%%%%%%%%

%%%%%%%%%%%%%%%%%%%%%%%%%%%%%%%%%%%%%%%%%%%%%
\begin{exo}[PCSI]{Enroulement d'un pendule}
%\begin{figure}[h!]
%\centering
\begin{center}
\includegraphics[width=0.9\columnwidth]{pendule.pdf}
\end{center}
%\end{figure}
Un point matériel de masse $m$ est suspendu à un point fixe A par un fil inextensible de longueur $l$ et de masse négligeable. On lâche ce pendule sans vitesse initiale, d'un angle $\alpha$ avec la verticale. Une tige horizontale fixe est placée à la verticale sous le point A à une distance $d<l$, de sorte que le fil heurte cette tige lorsque le pendule passe par sa position d'équilibre.  On admet que la vitesse du point matériel se conserve au cours du choc, et que le fil demeure tendu juste après le choc.
\begin{quest}
\item Avec le théorème de l'énergie mécanique, déterminer la vitesse du point matériel après le choc en un point M repéré par l'angle $\theta$ en fonction de $g,l,d,\theta,\alpha$. On considérera que l'énergie potentielle de pesanteur est nulle au point A $(z_A=0)$.
\item Montrer que la tension du fil en M vaut $T=mg\left[3\cos\theta+2\,\frac{d-l\cos\alpha}{l-d}\right].$
\item À quelle condition le fil reste-t-il tendu au cours du mouvement ?
\item Pour $\alpha=\pi/2$, déterminer la condition pour que le fil s'enroule autour de la tige tout en restant tendu.
\end{quest}
\end{exo}
%%%%%%%%%%%%%%%%%%%%%%%%%%%%%%%%%%%%%%%%%%%%%


%%%%%%%%%%%%%%%%%%%%%%%%%%%%%%%%%%%%%%%%%%%%%
\begin{exo}[PCSI]{Propulsion d'une bille de flipper}
  %\begin{figure}[h!]
  %\centering
  \begin{center}
  \includegraphics[width=\columnwidth]{flipper.pdf}
  \end{center}
  %\end{figure}
  
  Un flipper est modélisé par un plan incliné faisant un angle $\alpha$ avec l'horizontale. La bille, assimilée à un point matériel de masse $m$, est lancée à l'aide d'un ressort idéal de longueur à vide $l_0$ et de constante de raideur $k$. On s'intéresse au mouvement rectiligne de la bille lors de la propulsion. La bille est repérée le long de l'axe $(Ox)$ par son abscisse $x$ et l'origine du repère est prise à l'extrémité du ressort à vide.
  On néglige
  %dans un premier temps
  tout frottement et on donne $g=\SI{9.8}{m.s^{-2}}$, $\alpha=\SI{7}{\degree}$, $m=\SI{60}{g}$.
  \begin{quest}
  \item Le ressort se comprime d'une longueur $e=\SI{0.3}{mm}$ lorsque la bille, au repos, prend appui sur lui. En déduire la constante de raideur $k$ du ressort en fonction de $m,g,e$. Faire l'application numérique.
  \end{quest}
  On écarte la bille de sa position d'équilibre en comprimant le ressort d'une distance $a$. On lâche alors la bille sans vitesse initiale. Dans toute la suite, on étudie le mouvement de la bille tant que celle-ci est en contact avec le ressort.
  \begin{quest}
  \item Déterminer l’équation différentielle vérifiée par l’abscisse $x(t)$ de la bille.
  \item La résoudre en s’aidant des conditions initiales.
  \end{quest}
  \end{exo}
  %%%%%%%%%%%%%%%%%%%%%%%%%%%%%%%%%%%%%%%%%%%%%

%%%%%%%%%%%%%%%%%%%%%%%%%%%%%%%%%%%%%%%%%%%%%
\begin{exo}[PCSI]{Propulsion d'une boule de flipper}
%\begin{figure}[h!]
%\centering
\begin{center}
\includegraphics[width=\columnwidth]{flipper.pdf}
\end{center}
%\end{figure}

Un flipper est modélisé par un plan incliné faisant un angle $\alpha$ avec l'horizontale. La bille, assimilée à un point matériel de masse $m$, est lancée à l'aide d'un ressort idéal de longueur à vide $l_0$ et de constante de raideur $k$. On s'intéresse au mouvement rectiligne de la bille lors de la propulsion. La bille est repérée le long de l'axe $(Ox)$ par son abscisse $x$ et l'origine du repère est prise à l'extrémité du ressort à vide.
On néglige
%dans un premier temps
tout frottement et on donne $g=\SI{9.8}{m.s^{-2}}$, $\alpha=\SI{7}{\degree}$, $m=\SI{60}{g}$.
\begin{quest}
\item Le ressort se comprime d'une longueur $e=\SI{0.3}{mm}$ lorsque la bille, au repos, prend appui sur lui. En déduire la constante de raideur $k$ du ressort en fonction de $m,g,e$. Faire l'application numérique.
\end{quest}
On écarte la bille de sa position d'équilibre en comprimant le ressort d'une distance $a$. On lâche alors la bille sans vitesse initiale. Dans toute la suite, on étudie le mouvement de la bille tant que celle-ci est en contact avec le ressort.
\begin{quest}
\item Déterminer l'énergie potentielle totale de la bille en fonction de $x,k,m,g,\alpha$. On prendra pour origine de l'énergie potentielle le point O.
\item Par une approche énergétique, trouver une relation entre $\dot{x},x,k,m,g,\alpha,a,e$.
\item En déduire l'équation différentielle vérifiée par l'abscisse $x$ de la bille.
% On posera $\omega_0=\sqrt{k/m}$.
La résoudre et exprimer $x(t)$ en fonction de $a,e,\omega_0$.
\end{quest}
\end{exo}
%%%%%%%%%%%%%%%%%%%%%%%%%%%%%%%%%%%%%%%%%%%%%

%%%%%%%%%%%%%%%%%%%%%%%%%%%%%%%%%%%%%%%%%%%%%%%%%%%%%%%%%%%%
\begin{exo}[MPII][2][TD]{Chute d'un satellite Starlink}
  \Notions{\item Dynamique du point \task Gravitation \task Théorème énergétique\task Vitesse cosmique}
%\begin{wrapfigure}[6]{r}[34pt]{4.5cm}
%\vspace{-0.5cm}
%\includegraphics[scale=0.4]{star.jpeg}
%\end{wrapfigure}
Les satellites Starlink sont lancés régulièrement à des orbites basses, d'altitude $H=\SI{600}{km}$, dans le but de déployer un réseau internet mondial. Un satellite Starlink de masse $m$ est en orbite autour de la Terre de masse $M$. Le référentiel géocentrique est supposé galiléen. On note $G$ la constante de gravitation.

 \begin{quest}

%  \item On donne l'accélération en coordonnées polaires,
%\begin{equation}
% $\Vec{a} = (\ddot{r}-r\dot{\theta}^2) \er + (2\dot{r} \dot{\theta} + r \ddot{\theta}) \et.$
%\end{equation}
  % Donner l'expression de la force de gravitation qui s'exerce sur le satellite, puis projeter le PFD. 
  \item Montrer que $C=r^2\dot{\theta}$ est une constante à partir du théorème du moment cinétique ($2^{\rm e}$ loi de Képler). %Retrouver cette loi à partir du TMC.

  \sol{
  La loi de Newton de la gravitation donne
  $$
\Vec{F_G} = -\dfrac{GmM}{r^2}.
  $$
 Le PFD s'écrit alors
$$ m \Vec{a} = -\dfrac{GmM}{r^2} \er,
$$
ce qui donne après projection dans une base polaire $(\er,\et)$, les équations
$$
\left\{
\begin{array}{ll}
\ddot{r}-r\dot{\theta}^2 & =  -\dfrac{GM}{r^2},\\
2\dot{r} \dot{\theta} + r \ddot{\theta} & = 0.
\end{array}\right.
$$
En intégrant la $2^{\rm e}$ équation, on trouve  $2\dot{r}r \dot{\theta} + r^2 \ddot{\theta} = \dd{}{t}(r^2\dot{\theta})=0$, d'où l'on tire la $2^{\rm e}$ loi de Képler,

\eqbox{
C = r^2\dot{\theta} = \cte.
}

\udl{Remarque} : Pas besoin de supposer que $r=\cte$ pour démontrer la $2^{\rm e}$ loi de Képler !
  }

  \item On donne l'accélération en coordonnées polaires,
  %\begin{equation}
  $\Vec{a} = (\ddot{r}-r\dot{\theta}^2) \er + (2\dot{r} \dot{\theta} + r \ddot{\theta}) \et.$ Donner l'expression de la force de gravitation qui s'exerce sur le satellite, puis projeter le PFD. Dans l'hypothèse d'une trajectoire circulaire de rayon $r$, montrer  que la vitesse du satellite sur son orbite a pour expression
  $$
  v = \sqrt{\frac{GM}{r}}.
  $$
%Déterminer l'expression de la période $T$ et vérifier que $\frac{r^3}{T^2}$ est indépendant de la masse du satellite ($3^{\rm e}$ loi de Képler).
  \sol{
  On projette le PFD sur $\er$. Sachant que $v=r\dot{\theta}$, on trouve alors
  $$
  -r\dot{\theta}^2 = -\frac{v^2}{r}=  -\frac{GM}{r^2},
  $$
  soit

  \eqbox{
  v = \sqrt{\frac{GM}{r}}.
  }
 % La période s'obtient par
 %  $
 %  T = \frac{2\pi r}{v}  = 2\pi r  \sqrt{\frac{r}{GM}}.
 %  $
 %  On trouve alors
 %  $$
 %  \dfrac{r^3}{T^2} = \dfrac{GM}{4\pi^2}.
 %  $$
 }

  \end{quest}

  Sur cet orbite circulaire, on constate que l'altitude du satellite diminue à chaque période d'une hauteur $h \ll r$ en raison des frottements fluides avec l'atmosphère terrestre. On décrit ces frottements par une force quadratique $\Vec{F} = - \alpha m v \,\Vec{v}$. Le coefficient $\alpha$ est
une constante.

\begin{quest}

  \item Montrer que la variation d'énergie mécanique induite par les frottements au cours d'une période a pour expression
  $
  \Delta E_{\rm m} = -2\pi \alpha G m M.
  $

  \sol{
  En supposant l'effet des frottements « faible », la trajectoire reste quasi-circulaire et donc le déplacement élémentaire vaut $\dl = r\d\theta \et$.
  Le travail des forces de frottements sur une période vaut
  $$
  W = \int_{\text{orbite}} \Vec{F} \cdot \Vec{d \ell} = -\int_0^{2\pi} \alpha m v^2 r d \theta = -\int_0^{2\pi} \dfrac{\alpha G m M}{r} r d \theta =  -\alpha G m M \int d \theta = \Delta E_{\rm m}.
  $$
Le théorème de l'énergie mécanique donne la dernière égalité. Ansi,
\eqbox{
  \Delta E_{\rm m} = -2\pi \alpha G m M.
}
  }

  \item Déterminer l'expression de l'énergie mécanique $E_{\rm m}$. En déduire l'altitude $h$ perdue par période en fonction de $\alpha$, de l'altitude du satellite $H$ et du rayon terrestre $R_{\rm T}$. Application numérique.

  \sol{
  Il vient
  $$
  E_{\rm m} = E_{\rm c} + E_{\rm p}  = \dfrac{1}{2}mv^2 - \dfrac{GMm}{r} =  -\dfrac{GmM}{2r}.
  $$
    Au premier ordre en $h/r$, on trouve avec $h=r_2-r_1$,
  $
  \Delta E_{\rm m} = -\frac{GmM}{2}\left(\frac{1}{r_2}-\frac{1}{r_1}\right) \simeq -\frac{GmMh}{2r^2} =-2\pi \alpha G m M
  $
  et donc
\eqbox{
h = 4\pi r^2 \alpha,
}
  avec $r=H+R_{\rm T}$. L'application numérique donne $h=\SI{0.9}{m}$.
  }

%  \item En déduire l'expression de l'altitude $h$ perdue par période en fonction de $\alpha$, de l'atitude du satellite $H$ et du rayon terrestre $R_{\rm T}$. Application numérique.
%  %Au bout de combien de temps l'altitude aura-t-elle baissé de $\SI{10}{km}$ ?
%
%  \sol{
%  Au premier ordre en $h/r$, on trouve avec $h=r_2-r_1$,
%  $
%  \Delta E_{\rm m} = -\frac{GmM}{2}\left(\frac{1}{r_2}-\frac{1}{r_1}\right) = -\frac{GmMh}{2r^2} =-2\pi \alpha G m M
%  $
%  et donc $h = 4\pi r^2 \alpha$ avec $r=H+R_{\rm T}$. L'application numérique donne $h=\SI{}{}$.
%
%%  puis avec $\Delta h =\SI{10}{km}$,
%%  $$
%%  \Delta t = T \dfrac{\Delta h}{h} = \dfrac{\Delta h}{2\alpha \sqrt{rGM}} \simeq 2\,\text{ans}.
%%  $$
%  }


 \end{quest}
 \noindent
$\begin{array}{lll}
\textsc{Données :}
& \text{constante de gravitation} & :  G = \SI{6.67e-11}{m^3.kg^{-1}.s^{-2}} \\
& \text{coefficient de frottement} & : \alpha  = \SI{1.5e-15}{m^{-1}} \\
& \text{rayon terrestre} & : R_{\rm T}=\SI{6360}{km}
\end{array}
$

\end{exo}
%%%%%%%%%%%%%%%%%%%%%%%%%%%%%%%%%%%%


%%%%%%%%%%%%%%%%%%%%%%%%%%%%%%%%%%%%%%%%%%%%%%%%%%%%%%%%%%%%
\begin{exo}[MPII][2][devoir]{Chute d'un satellite sur Terre}
  \Notions{\item Dynamique du point \task Gravitation \task Théorème énergétique\task Vitesse cosmique}
%\begin{wrapfigure}[6]{r}[34pt]{4.5cm}
%\vspace{-0.5cm}
%\includegraphics[scale=0.4]{star.jpeg}
%\end{wrapfigure}
Les satellites Starlink sont lancés régulièrement à des orbites basses, d'altitude $H=\SI{600}{km}$, dans le but de déployer un réseau internet mondial. Un satellite Starlink est assimilé à un point matériel $M$ de masse $m$ est en orbite autour de la Terre de masse $M_T$. Le référentiel géocentrique est supposé galiléen. On note $G$ la constante de gravitation.

 \begin{quest}

 \item 
  Justifier en s'aidant du théorème du moment cinétique que le mouvement en plan. Calculer le moment cinétique $\Vec{L_O}$ du satellite par rapport au centre de la Terre $O$, et montrer que $C=r^2\dot{\theta}$ est une constante.
  %  ($2^{\rm e}$ loi de Képler).

  \sol{
  \udl{Référentiel} : référentiel géocentrique $\Rg$ supposé galiléen.\\ 

  \udl{Système} : satellite de centre $M$ et de masse $m$.\\ 

  \udl{Bilan des actions mécaniques} : la loi de Newton de la gravitation donne
  $$
\Vec{F_G} = -\dfrac{GmM_T}{OM^3}\Vec{OM}.
  $$
 Le moment de cette force, qui est centrale, est donc nulle, 
 $$
  \Vec{\mc{M}}(\Vec{F_G}) = \Vec{OM} \wedge \Vec{F_G} = \Vec{0}.
 $$
 \udl{Théorème du moment cinétique} :
 Le TMC indique donc que le moment cinétique est conservé, soit 
 \eqbox{
 \Vec{L_O} =  \Vec{Cte}.
 }
 Cela indique en particulier que le mouvement est compris dans un plan, que l'on suppose orthogonal à la direction $\ez$. En utilisant le coordonnées polaires $(r,\theta)$ dans ce plan, e moment cinétique du satellite a pour expression 
 $$
 \Vec{L_O} = \Vec{OM} \wedge m\Vec{v} = r\er \wedge m(\dot{r} \er + r\dot{\theta}\er) = mr^2\dot{\theta} \ez.
 $$
On en déduit que 
 \eqbox{
 C = r^2\dot{\theta} = \cte.
 }
% \udl{Remarque} : Pas besoin de supposer que $r=\cte$ pour démontrer la $2^{\rm e}$ loi de Képler !
  }

  \item On donne l'accélération en coordonnées polaires,
  $\Vec{a} = (\ddot{r}-r\dot{\theta}^2) \er + (2\dot{r} \dot{\theta} + r \ddot{\theta}) \et.$
  Dans l'hypothèse d'une trajectoire circulaire de rayon $r$, montrer que la vitesse du satellite sur son orbite a pour expression 
  % (première vitesse cosmique)
  \begin{equation}
    v = \sqrt{\frac{GM_T}{r}}.
  \end{equation}
%Déterminer l'expression de la période $T$ et vérifier que $\frac{r^3}{T^2}$ est indépendant de la masse du satellite ($3^{\rm e}$ loi de Képler).
  \sol{ 
  On applique cette fois-ci le PFD, que l'on projette dans la direction radiale $\er$.
  $$
  m \Vec{a} = -\dfrac{GmM_T}{r^2} \er \Rar \ddot{r}-r\dot{\theta}^2 =  -\dfrac{GM_T}{r^2}.
  $$
  Pour une trajectoire circulaire, on a $r=\cte$ et $v=r\dot{\theta}$, d'où
  $$
  -r\dot{\theta}^2 = -\frac{v^2}{r}=  -\frac{GM_T}{r^2}.
  $$
  On en tire l'expression de la première vitesse cosmique,
  \eqbox{
  v = \sqrt{\frac{GM_T}{r}}.
  }
 % La période s'obtient par
 %  $
 %  T = \frac{2\pi r}{v}  = 2\pi r  \sqrt{\frac{r}{GM}}.
 %  $
 %  On trouve alors
 %  $$
 %  \dfrac{r^3}{T^2} = \dfrac{GM}{4\pi^2}.
 %  $$
 }

  \end{quest}

  Sur cette orbite circulaire, on constate que l'altitude du satellite diminue à chaque période d'une hauteur $h \ll r$ en raison des frottements fluides avec l'atmosphère terrestre. On décrit ces frottements par une force quadratique $\Vec{F} = - \alpha m v \,\Vec{v}$. Le coefficient $\alpha$ est
une constante.

\begin{quest}

  \item Justifier que la force $\F$ est quadratique en la vitesse $v$. Montrer que la variation d'énergie mécanique induite par les frottements au cours d'une période a pour expression
  \begin{equation}\label{eq:deltaE}
    \Delta E_{\rm m} = -2\pi \alpha G m M_T.
  \end{equation}

  \sol{
    La vitesse du satellite est très élévée par rapport aux vitesses du quotidien, on peut donc supposer que les frottements fluides sont quadratiques en la vitesse et non linéaires. 
  En supposant l'effet des frottements « faible », la trajectoire reste quasi-circulaire et donc le déplacement élémentaire vaut $\dl = r\d\theta \et$.
  Le travail des forces de frottements sur une période vaut
  $$
  W = \int_{\text{orbite}} \Vec{F} \cdot \Vec{d \ell} = -\int_0^{2\pi} \alpha m v^2 r d \theta = -\int_0^{2\pi} \dfrac{\alpha G m M_T}{r} r d \theta =  -\alpha G m M_T \int d \theta = \Delta E_{\rm m}.
  $$
Le théorème de l'énergie mécanique donne la dernière égalité. Ansi,
\eqbox{
  \Delta E_{\rm m} = -2\pi \alpha G m M_T.
}
  }

  \item Déterminer l'expression de l'énergie mécanique $E_{\rm m}$ en fonction de la variable $r$. En déduire l'altitude $h$ perdue par période en fonction de $\alpha$, de l'altitude du satellite $H$ et du rayon terrestre $R_{\rm T}$. \AN.

  \sol{
  L'énegie mécanique $E_{\rm m}$ est la somme de l'énergie cinétique $E_{\rm c}$ et de l'énergie potentielle $E_{\rm p}$, soit
  $$
  E_{\rm m} = E_{\rm c} + E_{\rm p}  = \dfrac{1}{2}mv^2 - \dfrac{GmM_T}{r} =  -\dfrac{GmM_T}{2r}.
  $$
    % Au premier ordre en $h/r$, on trouve avec $h=r_2-r_1$,
    La variation d'énergie mécanique sur une période entre deux rayons $r_1=r$ et $r_2=r-h$ vaut donc
  $$
  \Delta E_{\rm m} = -\frac{GmM}{2}\left(\frac{1}{r_2}-\frac{1}{r_1}\right) = -\frac{GmM}{2r^2}\left(\frac{1}{1-h/r}-1\right) \simeq -\frac{GmMh}{2r^2},
  $$
  % au premier ordre en $h/r$.
  car $\frac{1}{1-x}  = 1+x +o(x)$ pour $x\ll 1$.
  D'après l'\Eq{deltaE}, on trouve donc
\eqbox{
h = 4\pi r^2 \alpha,
}
  où la distance vaut $r=H+R_{\rm T}$. L'application numérique donne \fbox{$h=\SI{0.9}{m}$}.
  }

%  \item En déduire l'expression de l'altitude $h$ perdue par période en fonction de $\alpha$, de l'atitude du satellite $H$ et du rayon terrestre $R_{\rm T}$. Application numérique.
%  %Au bout de combien de temps l'altitude aura-t-elle baissé de $\SI{10}{km}$ ?
%
%  \sol{
%  Au premier ordre en $h/r$, on trouve avec $h=r_2-r_1$,
%  $
%  \Delta E_{\rm m} = -\frac{GmM}{2}\left(\frac{1}{r_2}-\frac{1}{r_1}\right) = -\frac{GmMh}{2r^2} =-2\pi \alpha G m M
%  $
%  et donc $h = 4\pi r^2 \alpha$ avec $r=H+R_{\rm T}$. L'application numérique donne $h=\SI{}{}$.
%
%%  puis avec $\Delta h =\SI{10}{km}$,
%%  $$
%%  \Delta t = T \dfrac{\Delta h}{h} = \dfrac{\Delta h}{2\alpha \sqrt{rGM}} \simeq 2\,\text{ans}.
%%  $$
%  }


 \end{quest}
 \noindent
$\begin{array}{lll}
\textsc{Données :}
& \text{Constante de gravitation} & :  G = \SI{6.67e-11}{m^3.kg^{-1}.s^{-2}}, \\
& \text{Coefficient de frottement} & : \alpha  = \SI{1.5e-15}{m^{-1}}, \\
& \text{Rayon terrestre} & : R_{\rm T}=\SI{6360}{km}.
\end{array}
$

\end{exo}
%%%%%%%%%%%%%%%%%%%%%%%%%%%%%%%%%%%%

%%%%%%%%%%%%%%%%%%%%%%%%%%%%%%%%%%%%%%%
\begin{exo}[MPSI]{Masselotte immergée attachée à un ressort}

On considère un cylindre de masse $m$, de hauteur $h$ et de section $s$ relié à un ressort vertical de raideur $k$. Le cylindre  se déplace dans un champ de pesanteur $\Vec{g} = -g \Vec{e_z}$. Le ressort est fixé à un support immobile dans le référentiel du laboratoire supposé galiléen. On note $z$ l'altitude du centre de masse du cylindre, l'origine $z=0$ correspondant à la position d'équilibre. On ajoute ensuite de l'eau de masse volumique~$\rho$ jusqu'à une hauteur $z_{\rm s}$.
% ; le cylindre reste constamment vertical au cours du mouvement.

%%% FIGURE : %%%%
%\begin{figure}[h!]
%\centering
\begin{center}
\begin{tikzpicture}


\filldraw[draw opacity = 0, top color=gray!50!white, bottom color=white] (-2.25,0.5) -- ++ (5,0) -- ++ (0,-1) -- ++ (-5,0) -- cycle;
\draw (-2.25,0.5) -- ++ (5,0);
\fill[draw=black,fill=white] (0,0) rectangle (0.5,2);
\draw  plot[domain=0:10*pi, samples=100] ({0.25+0.1*sin(\x r)},2+0.05*\x);
\draw (-0.5,3.6) -- ++ (1.5,0);
\coordinate (O) at (0.25,1);
\node at (O) {$\bullet$};
\draw[->,>=latex] (-3,-1) -- ++ (0,5);
\draw[densely dashed] (O) -- ++ (-3.3,0) node[left] {$z$};
\draw[densely dashed] (-2.25,0.5) -- ++ (-0.75,0) node[left] {$z_{\rm s}$};

\draw[->,>=latex,JoliRouge] (2,3) -- ++ (0,-1) node[below] {$\vec{g}$}; 


\end{tikzpicture}
\end{center}
%\caption{}
%\label{fig:}
%\end{figure}
%%%%%%%%%%%%%


\begin{quest}
\item Donner l'expression des forces s'exercent sur le cylindre. On distinguera trois situations : celle où le cylindre est partiellement immergé, celle où le cylindre est complètement immergé, et celle où le solide est entièrement hors de l'eau. On négligera tout effet « dynamique » de l'eau sur le mouvement du cylindre.

\sol{
$\Vec{F_{\rm ressort}} = -k(z-\ell) \Vec{e_z}$, $\Vec{P} = -mg\Vec{e_z}$ et à l'équilibre $-k(0-\ell)-mg = 0$ et donc $\ell =  mg/k$. Poussée d'Archimède $\Vec{\Pi} = -\rho (z-z_s-h/2) s g\Vec{e_z}$ si $z \in [z_s-h/2,z_s+h/2]$ (partiellement immergé), $\Vec{\Pi} = \rho h s g\Vec{e_z}$ si $z < z_s - h/2$ (complètememt immergé), et $\Vec{\Pi} = \Vec{0}$ si $z>z_s + h/2$.

}
\item Appliquer le principe fondamental de la dynamique au centre de masse du cylindre dans ces trois situations. Montrer que lorsque le cylindre est partiellement immergé, la dynamique est celle d'une masselotte reliée à un ressort de raideur plus élevée $k'$ dont on donnera l'expression et la valeur numérique.


\sol{
Le PFD donne, lorsque $z \in [z_s-h/2,z_s+h/2]$,
$$
m\ddot{z} = -k(z-\ell) - mg -  s\rho g(z-z_{\rm s}-h/2) = -(k+\rho s g)z +   s\rho g(z_{\rm s}+h/2).
$$
Raideur $k' = k +\rho s g$.
ou lorsque $z< z_s - h/2$,
$$
m\ddot{z} = -kz + s\rho gh.
$$
ou lorsque $z> z_s + h/2$,
$$
m\ddot{z} = -kz.
$$

}

\item Déterminer la position d'équilibre en présence d'eau $z_0$. Que doit valoir le niveau d'eau $z_s$ pour que cette position d'équilibre en présence d'eau coincide avec l'origine ? Commenter.

\sol{
Selon les cas, on trouve
$$
z_0 = \dfrac{\rho s g}{k'} (z_s + h/2), \qquad z_0 = \dfrac{\rho s g }{k'} h, \qquad z_0=0.
$$
On a $z_0=0$ pour $z_s < -h/2$. Le cylindre est à l'air libre.
}

\item  On suppose que $z_{\rm s} =0$. Quelle est l'équation différentielle vérifiée par $z(t)$ en fonction de la situation (cylindre hors de l'eau, partiellement immergé, complètement immergé) ? On introduira deux grandeurs $\omega_0$ et $\omega_0'$ homogènes à une pulsation.
Applications numériques.

\sol{
On a
$\ddot{z} + \omega^2 z = \omega^2 z_0$ avec  $\omega = \omega_0$ ou $\omega_0'$ (partiellement immergé), et $z_0 = 0$ ou $\rho s g h /k$ (complètement immergé).
$$
\omega_0 = \sqrt{\dfrac{k}{m}},  \qquad \omega_0' = \sqrt{\dfrac{k'}{m}} > \omega_0.
$$
}

\item On lâche le cylindre sans vitesse initiale d'une hauteur $z=h/2$. Déterminer la trajectoire du centre de masse du cylindre $z(t)$ jusqu'au point le plus bas $z_{\rm min}$, que l'on déterminera. Applications numériques. Représenter graphiquement l'allure de $z(t)$ sur plusieurs périodes.

\sol{
$z > 0$ donc le cylindre est d'abord hors de l'eau donc l'équation à résoudre est
$$
\ddot{z} + \omega_0^2 z  = 0
$$
de solution $z(t) = h/2 \cos(\omega_0 t)$. Le cylindre s'immerge partiellement lorsque $z=0$, soit au bout d'un temps
$t_1 = \dfrac{\pi}{2\omega_0}$.
La vitesse vaut alors $\dot{z}(t_1) = h \omega_0 /2$. Pour $t>t_1$, l'équation à résoudre est
$$
\ddot{z} + \omega_0'^2 z  = 0
$$
de solution $z(t) = A \cos(\omega_0' t) + B \sin(\omega_0' t)$ avec (on ramène l'origine  à $t_1$,
$$
z(0) = A = 0
$$
et
$$
\dot{z}(t_1) = B \omega_0' = h \omega_0/2 \Rightarrow B = \dfrac{h}{2} \sqrt{\dfrac{k}{k'}} < \dfrac{h}{2}.
$$
Cette équation horaire est vérifiée jusqu'à ce que $z=0$ de nouveau. Le minimum est atteint losque le sinus vaut $-1$, soit

$$
z_{\min} = - \dfrac{h}{2} \sqrt{\dfrac{k}{k'}}.
$$
}

\end{quest}

$\begin{array}{lll}
\textsc{Données :}  & \text{masse volumique de l'eau} & : \rho=\SI{1.0e3}{kg.m^{-3}} \\
& \text{intensité de la pesanteur} & : g=\SI{9.8}{m.s^{-2}} \\
& \text{raideur du ressort} & : k=\SI{20}{N.m^{-1}} \\
& \text{rayon du cylindre } & : r=\SI{1.0}{cm} \\
& \text{hauteur du cylindre} & : h=\SI{10}{cm} \\
& \text{masse du cylindre} & : m=\SI{100}{g}
\end{array}
$
\end{exo}
%%%%%%%%%%%%%%%%%%%%%%%%%%%%%%%%%%%%%%%%%%

%%%%%%%%%%%%%%%%%%%%%%%%%%%%%%%%%%%%%%
\begin{exo}[MPSI]{Atterrissage d'un avion de chasse}
 Un avion de chasse de masse $m=\SI{9}{t}$, ayant les freins en panne, aborde une piste d'atterrissage avec une vitesse $v_{0}=\SI{241}{km.h^{-1}}$. Au sol,
 il est freiné par un parachute de diamètre $D=3\:\mathrm{m}$ déployé instantanément au contact du sol. La force de traînée
 du parachute est $\Vec{T} = - \textstyle\frac{1}{2}C\rho Sv \Vec{v}$ avec $\Vec{v}=v\ex$ la vitesse de l'avion, $S$ la surface du parachute projetée sur le plan perpendiculaire à la direction du mouvement, $C=\SI{1.5}{}$ le
 coefficient de traînée, et $\rho=\SI{1.2}{kg.m^{-3}}$ la masse volumique de l'air.
 \begin{quest}
  \item  On néglige les forces de frottement solide.
	\begin{enumerate}
	 \item Déterminer l'une équation différentielle vérifiée par $\Vec{v}$ et l'intégrer.

	 \sol{
   \begin{itemize}
    \item \udl{Référentiel} : référentiel terrestre $\RT$ supposé galiléen.
    \item \udl{Système} : $\{$avion$\}$, assimilé à un point matériel $M$ de masse $m$.
    \item \udl{Bilan des actions mécaniques} : 
    \begin{enumerate}
      \item Force de traînée du parachute $\Vec{T} = - \frac{1}{2}C\rho Sv^2 \ex $ avec $S=\pi D^2/4$.
      \item Force de pesanteur $\Vec{P} = -mg\Vec{e_z}$.
      \item Force de réaction du sol $\Vec{R} = R \ez$.
    \end{enumerate}
    
    \item \udl{PFD} :
    $$
    m\dd{\Vec{v}}{t} = \Vec{T} + \Vec{P} + \Vec{R} \Rightarrow \left\{
    \begin{array}{ll}
    m\dot{v} = -\dfrac{C\rho S}{2}v^{2},\\[3mm]
    0 = -mg +R.
    \end{array}
    \right.
    $$
    La projection sur l'axe horizontal $(Ox)$ donne
    $$\dot{v}=-\frac{C\rho S}{2m}v^{2}.$$
    On voit apparaître une distance caractéristique de freinage
    \eqbox{
      k = \frac{2m}{C\rho S} = \SI{540}{m}.
      }
    L'équation différentielle vérifiée par la vitesse est donc
    \eqbox{
      \dd{v}{t} = -\frac{v^{2}}{k}.
    }
  \end{itemize} 
  On résout cette équation différentielle par séparation des variables :
          $$
           \frac{\d v}{v^2} = -\d\left( \frac{1}{v} \right) = \frac{\d t}{k} \Rar \int_{v_0}^{v(t)} \d\left( \frac{1}{v} \right) = \frac{1}{v(t)} - \frac{1}{v_0} = -\frac{t}{k}.
          $$
	   d'où l'on tire 
     \eqbox{
     v(t) = \frac{v_0}{1 + \frac{v_0}{k}t}.
     }
     Le temps typique de freinage dépend de la vitesse initiale et vaut 
     \eqbox{
     \tau = \frac{k}{v_0} = \SI{8.1}{s}.
     }
	 }

	 \item Montrer que la force de traînée n'est pas suffisante pour arrêter l'avion.

	 \sol{
    Pour arrêter l'avion, il faudrait que la vitesse s'annule. Or
	 $v = 0$ lorsque $t \rightarrow +\infty$. Donc l'avion se s'arrête pas en un temps fini.
	 }

	\end{enumerate}
  \item Avec des freins, l'avion subit en plus une force de frottement solide qui lui permet de s'arrêter à temps.
	\begin{enumerate}

\item
  La distance d'arrêt vaut $d=\SI{1,4}{km}$. Déterminer la vitesse qu'aurait l'avion à cette distance s'il était freiné uniquement par le parachute, donc sans la force de frottement solide.
  \sol{
    Deux méthodes sont possibles.
    \begin{enumerate}
      \item On intègre $v(t)$ pour obtenir $x(t)$, on déterminer l'instant $t_d$ tel que $x(t_d)=d$, puis on évalue $v(t_d)$.
      \item On détermine l'équation vérifiée par la fonction $v(x)$, puis on résout et on évalue $v(x=d)$.
    \end{enumerate}
Appliquons ces deux méthodes.
    \begin{enumerate}
      \item 
      La loi horaire $x(t)$ s'obtient par intégration de la loi horaire $v(t)$, soit
      $$
     x(t) = \int_0^t v(t')\d t' = \int_0^t \frac{v_0}{1 + \frac{v_0}{k}t'} \d t' = 
      k\ln\left(1 + \frac{v_0}{k}t\right).
       $$
      On inverse cette relation pour déterminer l'instant d'arrêt $t_d$,
      $$
      d = k\ln\left(1 + \frac{v_0}{k}t_d\right) \Rar t_d = \frac{k}{v_0}\left(e^{d/k} - 1\right).
      $$
      En en déduit la vitesse qu'aurait l'avion à ce instant s'il était freiné uniquement par le parachute,
      \eqbox{
      v(t_d) = \frac{v_0}{1 + \frac{v_0}{k}t_d} = v_0\, e^{-d/k} = \SI{18}{km.h^{-1}}.
      }
      \item L'équation vérifiée par la vitesse se réarrange en
      $$
      \dd{v}{t} = \dd{v}{x}\dd{x}{t} = \dd{v}{x} v = -\frac{v^2}{k} \Rar \dd{v}{x} = -\frac{v}{k}.
      $$
      On obtient une équation différentielle linéaire du premier ordre, qui se résout directement en 
      $$
      v(x) = v_0 \, e^{-{x}/{k}}.
      $$
      On en déduit ensuite $v(d)$.
    \end{enumerate}
  }
  \item On modélise la force de frottement solide par une force tangentielle opposée au sens du mouvement de norme constante $F = \SI{5e3}{N}$. Montrer que la vitesse vérifie l'équation différentielle
  \begin{equation}
  \dfrac{\dot{v}}{v^{2} + w^2}=-\frac{1}{k},
  \end{equation}
  où $w$ et $k$ sont des constantes que l'on exprimera et calculera. Combien de temps met l'avion pour s'arrêter d'après ce modèle ? 
  
    \sol{
      On reprend le PFD de la première questions en incluant la force $-F\ex$. On trouve sur l'axe $(Ox)$,
      $$
    m \dot{v} = -\frac{C\rho S}{2}v^{2} - F \Rar \dot{v} = -\frac{C\rho S}{2m}v^{2} - \frac{F}{m} = -\frac{v^{2}}{k} - \frac{F}{m} = -\frac{1}{k}\left(v^{2} + w^2\right),
      $$
      où la vitesse caractéristique associée au frottement solide $w$ a pour expression
  \eqbox{
    w = \sqrt{\dfrac{kF}{m}} = \SI{17}{m.s^{-1}}.
    }
  On obtient bien l'équation différentielle demandée,
  $$
  \dfrac{\dot{v}}{v^{2} + w^2}=-\frac{1}{k}.
  $$
  On intègre cette équation par séparation des variables entre $t=0$ et le temps d'arrêt $\tau$,
  $$
  \frac{\d v}{v^2 + w^2} = -\frac{\d t}{k} \Rar \int_{v_0}^{0} \frac{\d v}{v^2 + w^2} = -\frac{1}{k} \int_0^t \d t \Rar 
  -\dfrac{1}{w} \arctan \left(\dfrac{v_0}{w}\right) = -\dfrac{\tau}{k},
  $$
  soit,
  \eqbox{
  \tau = \frac{k}{w}\arctan \left(\dfrac{v_0}{w}\right) \simeq \frac{\pi k}{2 w},
  }
  dans la limite $v_0 \gg w$, qui est vérifiée dans notre cas. On trouve numériquement,
  \eqbox{
  \tau  \simeq \SI{50}{s}.
  }
  }

  \end{enumerate}
 \end{quest}

 \textsc{Données :} une primitive de $v \mapsto \displaystyle \dfrac{1}{v^2+w^2}$ est $\displaystyle\dfrac{1}{w}\arctan\left(\dfrac{v}{w}\right)$.

\end{exo}
%%%%%%%%%%%%%%%%%%%%%%%%%%%%%%%%%%%%%%%%%%


%%%%%%%%%%%%%%%%%%%%%%%%%%%%%%%%%%%%%%
\begin{exo}[MPSI]{Atterrissage d'un avion}
 Un avion de chasse de masse $m=\SI{9}{t}$ aborde une piste d'atterrissage avec une vitesse $v_{0}=\SI{241}{km.h^{-1}}$. Au sol,
 il est freiné par un parachute déployé instantanément au contact du sol. La force de trainée
 du parachute vaut $\Vec{F} = - \lambda \Vec{v}$ avec $\Vec{v}=v\Vec{e_x}$ la vitesse de l'avion et $\lambda = \SI{150}{N.m^{-1}.s}$ le coefficient de frottement.
 \begin{quest}
  \item  On néglige les forces de frottement solide.
	\begin{enumerate}
	 \item Donner une équation différentielle vérifiée par $\Vec{v}$ et l'intégrer.

	 \sol{On applique le PFD selon l'horizontale : $\displaystyle\dot{v}=-\lambda v$. Ainsi
 \begin{equation}
	         v(t) = v_0 e^{-\lambda t}, \label{eq:avion1}
	        \end{equation}
	   où $v_{0}$ est la vitesse de l'avion au contact avec la piste d'atterrissage.
	 }

	 \item Montrer que la force de trainée n'est pas suffisante pour stopper l'avion.

	 \sol{
	 $v \underset{t\rightarrow+\infty}{\longrightarrow}0$.
	 }

	\end{enumerate}
  \item Avec des freins, l'avion subit en plus une force de frottement solide qui lui permet de s'arrêter à temps.
	\begin{enumerate}

\item
  La distance de freinage avec frottement solide est $d=\SI{1,4}{km}$. Déterminer la vitesse qu'aurait l'avion à cette distance s'il était freiné uniquement par le parachute.
  \sol{
  En intégrant l'équation~\eqref{eq:avion1}, on obtient
	  \begin{equation}
           x = \frac{v_0}{\lambda}(1-e^{-\lambda t}) \Rar d = \frac{v_0}{\lambda}(1-e^{-\lambda \tau})
          \end{equation}
	  donc \begin{equation}
	        v = v_{0}\,e^{-\lambda \tau} = v_0 \left(1- \frac{\lambda d}{v_0} \right) = v_0 - \lambda d.
	       \end{equation}
	  Ainsi, en $d$, $v=\SI{}{m.s^{-1}}$.
  }
  \item On modélise la force de frottement solide par une force tangentielle de norme $F = \SI{5e3}{N}$ constante et opposée au sens du mouvement. Déterminer l'équation différentielle
vérifié par la vitesse. Combien de temps met l'avion pour s'arrêter ?

    \sol{PFD puis
  $$
  w = \sqrt{\dfrac{2F}{C\rho S}}.
  $$
  On obtient
  $$
  \dfrac{1}{w} \left[ \arctan \left(\dfrac{v}{w}\right) - \arctan \left(\dfrac{v_0}{w}\right) \right] = -\dfrac{t}{w\tau}
  $$
  avec $$
  \tau = \dfrac{2m}{C\rho S w} = 2m \sqrt{2F C \rho S}.
  $$
  Donc
  $$
  v(t) = w \tan \left(\arctan \left(\dfrac{v_0}{w}\right)-\dfrac{t}{\tau} \right).
  $$
  La vitesse s'annule lorsque
  $t = \tau \arctan \left(\dfrac{v_0}{w}\right)	$.
  }

  \end{enumerate}
 \end{quest}

 % \textsc{Données :} une primitive de $v \mapsto \displaystyle \dfrac{1}{v^2+w^2}$ est $\displaystyle\dfrac{1}{w}\arctan\left(\dfrac{v}{w}\right)$.

\end{exo}
%%%%%%%%%%%%%%%%%%%%%%%%%%%%%%%%%%%%%%%%%%

%%%%%%%%%%%%%%%%%%%%%%%%%%%%%%%%%%%%
\begin{exo}[MPSI]{Déviation d’un faisceau d’électrons}
Un faisceau d'électron entre dans un condensateur plan de différence de potentiel $U$ avec une vitesse initiale $\Vec{v_{0}} = v_{0}\,\Vec{e_{x}}$. La distance entre les
 plaques est notée $d$, la longueur du condensateur est notée $l$. On néglige les effets de bord du condensateur. Ce type de dispositif est utilisé dans les tubes cathodiques des anciens postes de télévision.

  \begin{center}
 \includegraphics[scale=0.8]{Develec.eps}
\end{center}

 \begin{quest}
   \item Calculer la trajectoire des électrons.

   \sol{

   Le champ électrique entre les plaques du condensateur est $\Vec{E}=-U/d \,\Vec{e_{y}}$. Il soumet les électrons de charge $-e$ à une force $\Vec{f} = -e\Vec{E}$,
  donc
    \begin{equation}
     \begin{cases}
      v_{x} = v_{0}\\
      v_{y} = \dfrac{eU}{md} \,t
     \end{cases}.
    \end{equation}
    La trajectoire des particules est donc $y = \dfrac{eU}{2mdv_{0}^{2}}\,x^{2}$.
    }

  \item Calculer l'angle $\beta$ de déviation du faisceau.

  \sol{

 $\tan(\beta) = \dfrac{v_{y}(x=l)}{v_{x}(x=l)}=\dfrac{eUl}{mdv_{0}^{2}}$.

 }

  \item On place un écran à une distance $L$ de la sortie du condensateur. Montrer que la coordonnée $y$ du point d'impact des électrons sur l'écran est proportionnelle
  à $U$.

  \sol{

  Hors du condensateur, les charges ne sont soumises à aucune force, elles se déplacent en mouvement rectiligne uniforme avec les positions et vitesses en sortie
  du condensateur comme conditions initiales, donc $y(x=L)=\dfrac{el(L-l/2)}{mdv_{0}^{2}}\,U$.
  }

  \end{quest}


\end{exo}
%%%%%%%%%%%%%%%%%%%%%%%%%%%%%%%%%%%%%%%%

%%%%%%%%%%%%%%%%%%%%%%%%%%%%%%%%%%%%%%%%%
\begin{exo}[MPSI]{Mouvement dans un cyclotron}
Le cyclotron est un type d’accélérateur de particules inventé au début des années 1930. Il se compose de deux demi-disques dans lesquels règnent les champ magnétiques uniformes $\Vec{B}$ et $-\Vec{B}$. Entre ces deux demi-disques règne un champ électrique uniforme $\Vec{E}$. On considère une particule de charge $q$ et de masse $m$ en mouvement dans ce cyclotron.
  \begin{center}
 \includegraphics[scale=0.8]{Cyclotron.eps}
\end{center}
 \begin{quest}
  \item Montrer que le mouvement de la particule dans un des demi-disques est circulaire et uniforme. Exprimer le rayon $r$ de la trajectoire en fonction du champ magnétique $B$. Déterminer le temps $t_B$ passé par la particule dans chacun des demi-disques.

  \sol{
	$r=\dfrac{mv}{qB}$ et $t_1 = \dfrac{\pi r}{v} = \dfrac{\pi m}{qB}$.
  }


  \item Montrer que le mouvement de la particule entre les demi-disques est rectiligne accéléré. Soit $L$ la distance entre les faces planes des demi-disques. Quel est le gain de vitesse $\Delta v$ aquis par la particule entre les demi-disques ?

  \sol{
  $x = \dfrac{qE}{2m} t_2^2$ donc
  $$\Delta v = \dfrac{qE}{m} t_2 = \sqrt{\dfrac{2qEL}{m}}$$.
  }
  \item On suppose que $L$ est suffisamment petit pour le temps passé entre les demi-disques soi négligeable. Le champ électrique est maintenant supposé sinusoïdal, $E(t) = E\cos(2\pi f t)$. Quelle est la fréquence $f$ à donner au champ électrique pour que l'accélération de la particule soit optimale ?

  \sol{
  Fréquence cyclotron $f = \dfrac{qB}{2\pi m}$
  }

  \item La particule a une vitesse initiale nulle et part du bord d’un demi-disque. Combien faut-il faire de tour pour atteindre une vitesse $v=25 \times 10^{3}\ \mathrm{km.s^{-1}}$ pour un champ magnétique $B = 1$ T et une tension $U=\SI{2.5e3}{V}$ ?

  \sol{
  $2 N \Delta v = v$ donne $N$.
  }

\item Soit $R$ le rayon de l’armature qui entoure les demi-disques. Exprimer l’énergie cinétique maximale $E_c$ atteinte par la particule. Comment choisir $R$ et $B$ pour améliorer le cyclotron.

\sol{
$E_c = \dfrac{q^2B^2R^2}{2m}$
}
 \end{quest}


\end{exo}
%%%%%%%%%%%%%%%%%%%%%%%%%%%%%%%%%%%%%%%%%%%%%

%%%%%%%%%%%%%%%%%%%%%%%%%%%%%%%%%%%
\begin{exo}[PC][1][colle]{Force magnétique}
  
  Une bille d'acier de masse \(m=200 \mathrm{~g}\), fixée à l'extrémité d'un fil de longueur \(\ell=50 \mathrm{~cm}\) est attirée par un aimant de telle sorte qu'à l'équilibre, le fil s'incline de \(\alpha=50^{\circ}\). On prendra le champ de pesanteur terrestre \(g=9,8 \mathrm{~m} \cdot \mathrm{s}^{-2}\).
  
  Calculer la force magnétique que ressent la bille d'acier ainsi que la tension du fil.
  
  \begin{center}
    \includegraphics[width=0.5\textwidth]{2024_01_18_89e1c551bd3f28b92a09g-1(1)}
  \end{center}
  
\end{exo}
%%%%%%%%%%%%%%%%%%%%%%%%%%%%%%%%%%%

%%%%%%%%%%%%%%%%%%%%%%%%%%%%%%%%%%%
\begin{exo}[PC][1][colle]{Deux ressorts en équilibre}
  
  Considérons deux ressorts identiques, de même constante de raideur \(k=20 \mathrm{~N} \cdot \mathrm{m}^{-1}\), de longueur à vide \(\ell_{0}=20 \mathrm{~cm}\) montés comme l'indique la figure et soutenant deux points matériels \(\mathbf{M}_{1}\) et \(\mathbf{M}_{2}\) de masse \(m_{1}=100 \mathrm{~g}\) et \(m_{2}=150 \mathrm{~g}\). On prendra le champ de pesanteur \(g=10 \mathrm{~N} \cdot \mathrm{kg}^{-1}\).
  
  \begin{quest}
    \item En étudiant l'équilibre de \(\mathbf{M}_{2}\), calculer la longueur \(\ell_{2}\) du deuxième ressort.
    
    \item En étudiant l'équilibre de \(\mathbf{M}_{1}\), calculez la longueur \(\ell_{1}\) du premier ressort.
    
  \end{quest}
  
  \begin{center}
    \includegraphics[width=0.5\columnwidth]{2024_01_18_89e1c551bd3f28b92a09g-1(2)}
  \end{center}
\end{exo}
%%%%%%%%%%%%%%%%%%%%%%%%%%%%%%%%%%%

%%%%%%%%%%%%%%%%%%%%%%%%%%%%%%%%%%%
\begin{exo}[PC][1][colle]{Deux blocs en mouvement}
  
  Deux blocs de masse \(m\) (bloc A) et \(2 m\) (bloc B) sont mis en mouvement par une force extérieure horizontale \(\Vec{F}\).
  
  \begin{center}
    \includegraphics[width=0.5\textwidth]{2024_01_18_89e1c551bd3f28b92a09g-1}
  \end{center}
  
  \begin{quest}
    \item On suppose que le support sur lequel glissent les blocs ne produit aucun frottement. Exprimer l'action \(\Vec{F^{\prime}}\) du bloc \(\mathrm{A}\) sur le bloc \(\mathrm{B}\) en fonction des données du problème.
    
    \item Répondre à la même question en supposant maintenant que le contact entre les blocs et le support est caractérisé par un coefficient de frottement \(\mu\).
    
  \end{quest}
\end{exo}
%%%%%%%%%%%%%%%%%%%%%%%%%%%%%%%%%%%

%%%%%%%%%%%%%%%%%%%%%%%%%%%%%%%%%%%
\begin{exo}[PCSI][1][colle]{Mise en orbite d'une sonde spatiale}
On souhaite mettre en orbite une sonde spatiale $(s)$ autour de Mars. Ce véhicule possède au point $A$ la vitesse $\Vec{v_A}$ et présente un \guill{paramètre d'impact} $b$.


\begin{minipage}{0.5\columnwidth}
  \textsc{Données} :
  \begin{liste}
    \item
     $v_A=2,5 \mathrm{~km} \cdot \mathrm{s}^{-1}$;
    \item rayon de Mars : $r_M=3397 \mathrm{~km}$;
    \item masse de Mars : $M_M=6,39 \cdot 10^{23} \mathrm{~kg}$.
  \end{liste}
\end{minipage}\hfill
\begin{minipage}{0.45\columnwidth}
  % \begin{fig}[label]
    % \centering
    \includegraphics[width=\columnwidth]{orbite_mars.png}
    %\capt{}
  % \end{fig}
\end{minipage}
\begin{quest}
 \item Au point $A$, on suppose que le véhicule est assez éloigné de Mars pour pouvoir négliger l'énergie gravitationnelle. En déduire l'expression de l'énergie mécanique et la nature de la trajectoire de $(S)$. Montrer que le moment cinétique se conserve et définir la constante des aires $C$. Calculer $C$ en fonction de $v_A$ et $b$.
  Sachant que la trajectoire d'approche est tangente au cercle de rayon $\alpha \cdot r_M$ en $B$, calculer $v_B$ (vitesse en C) en fonction de $v_A, r_M, \alpha$ et $b$.
 \item Exprimer le paramètre d'impact $b$, en fonction de $r_M, v_A, M_M$ et $\alpha$. Application numérique pour $\alpha=3$.
 \item Déterminer la distance minimale $b_m$ pour que le véhicule évite la surface de Mars.
 \item Déterminer la vitesse $v_c$ d'un objet sur l'orbite circulaire de rayon $3 r_M$ ainsi que sa période de révolution en fonction des données. Effectuer l'application numérique.
 \item Au point $C$ (avec $\alpha=3$ ), on veut que le véhicule passe sur l'orbite circulaire de rayon $3 r_M$. Déterminer la variation de vitesse $\Delta v_c$ à communiquer au véhicule en fonction des données. Application numérique.
\end{quest}
\end{exo}
%%%%%%%%%%%%%%%%%%%%%%%%%%%%%%%%%%%

%%%%%%%%%%%%%%%%%%%%%%%%%%%%%%%%%%%
\begin{exo}[PC][1][colle]{Frottement sur un toboggan}
  Un enfant - que l'on assimilera à un point matériel $\rm M$ de masse $m=40 \mathrm{~kg}$ - glisse sur un toboggan décrivant une trajectoire circulaire de rayon $r=2,5 \mathrm{~m}$. L'enfant, initialement au point $\mathrm{A}$, d’angle $\theta=0^{\circ}$, se laisse glisser avec une vitesse initiale nulle, et atteint le point $\mathrm{B}$, d’angle $\theta=\SI{90}{^\circ}$ avec une vitesse $v_{\max}$. On supposera le référentiel terrestre  $\Rcal$ galiléen. 
  % On donne $g=9,8 \mathrm{~m} \cdot \mathrm{s}^{-2}$.

%%%%%%%%%%%%%%%%%%%%%%%%%%%%%%%%%%%%%%%%%%%
\begin{fig}[label]
\centering
\includegraphics[width=0.5\columnwidth]{2024_02_05_ce2a243e69fdb9afdd48g-1}
\capt{}
\end{fig}
%%%%%%%%%%%%%%%%%%%%%%%%%%%%%%%%%%%%%%%

\begin{quest}
  \item On suppose, dans un premier temps, que les frottements sont négligeables. Faire un bilan et un schéma des forces qui s'exercent sur $\rm M$.

  \item  À l'aide du principe fondamental de la dynamique,  exprimer la vitesse $v$ du point $\rm M$ en fonction de $\theta, r$ et le champ de pesanteur $g$. Exprimer la force qu'exerce le toboggan sur l'enfant en fonction de $m$, $g$ et $\theta$.


  \item Calculer la vitesse $v_{\max }$ acquise en $\mathrm{B}$.

\end{quest}

En réalité, les frottements induits par le contact du toboggan ne sont pas négligeables. Ils sont caractérisés par un coefficient de frottement $\mu=\num{0.2}$.

\begin{quest}
  \item Montrer que la force de frottement s'écrit approximativement
$
f \simeq 3 \mu m g \sin \theta.
$
 Expliquer en quoi il s'agit d'une approximation.

  \item Exprimer le travail produit par la force de frottement entre les points A et B. En déduire une nouvelle expression de la vitesse $v_{\max }$. Faire l'application numérique. Commenter.

\end{quest}
\end{exo}
%%%%%%%%%%%%%%%%%%%%%%%%%%%%%%%%%%%%%%%%%%%%


%%%%%%%%%%%%%%%%%%%%%%%%%%%%%%%%%%%
\begin{exo}[PC][1][colle]{Glissement d'un anneau}
Un petit anneau se déplace sur un guide qui a la forme d'une hélice. Les équations en coordonnées cylindriques de cette hélice rigide d'axe vertical ascendant $(\mathrm{O} z)$ sont

$$
\left\{\begin{aligned}
r & =a \\
z & =p \frac{\theta}{2 \pi}
\end{aligned}\right.
$$

où $a$ est le rayon de l'hélice et $p$ le pas de l'hélice.

L'anneau est abandonné sans vitesse initiale à partir d'une hauteur $z=h$. On assimile l'anneau à un point matériel glissant sans frottement le long de l'hélice.

\begin{quest}
  \item Relier la vitesse $v$ à $z$ au cours du mouvement, à l'aide du théorème de l'énergie mécanique.
\end{quest}
\end{exo}
%%%%%%%%%%%%%%%%%%%%%%%%%%%%%%%%%%%%%%%%%%%%


%%%%%%%%%%%%%%%%%%%%%%%%%%%%%%%%%%%
\begin{exo}[PC][1][colle]{Équilibre d'une barre}
On considère une barre maintenue horizontalement contre un mur à l'aide d'un fil fixé à son extrémité et à un point du mur (cf. figure ci-dessous). La barre a pour masse $m$ et pour longueur $\ell$, le fil quant à lui est de masse négligeable et fait un angle $\alpha=70^{\circ}$ par rapport au mur. On note $\Vec{T}$ l'action du fil sur la barre, $\Vec{P}$ le poids de la barre, $\Vec{R}$ l'action du mur sur la barre et $\Vec{g}$ le champ de pesanteur.

%%%%%%%%%%%%%%%%%%%%%%%%%%%%%%%%%%%%%%%%%%%
\begin{fig}[label]
\centering
\includegraphics[width=0.7\columnwidth]{2024_03_02_a2775ce5ed5c2ca3278ag-2}
% \capt{}
\end{fig}
%%%%%%%%%%%%%%%%%%%%%%%%%%%%%%%%%%%%%%%

\begin{quest}
  \item La droite d'action du poids coupe le fil au point $\mathrm{C}$. Calculer le moment du poids et de la tension au point $C$. En déduire que les trois forces $\Vec{T}, \Vec{P}$ et $\Vec{R}$ son concourrantes au point $C$, c'est-à-dire que ces trois forces ont leur droite d'action qui passe par le point $C$.

  \item Exprimer $R$ et $T$ en fonction de $m, g$ et $\alpha$.

  \item Le contact entre le mur et la barre est caractérisé par un coefficient de frottement $\mu=0,4$. Montrer que la barre ne glisse pas.

  \item On remplace le fil par un fil plus long de sorte que l'angle $\alpha$ diminue de $20 \%$. La barre restera-telle en équilibre?

\end{quest}
\end{exo}
%%%%%%%%%%%%%%%%%%%%%%%%%%%%%%%%%%%%%%%%%%%%

%%%%%%%%%%%%%%%%%%%%%%%%%%%%%%%%%%%
\begin{exo}[PC][1][oral]{Remplacement d'un ressort}
On souhaite remplacer le ressort qui permet de lancer une bille en acier à partir d'une rampe de lancement légèrement inclinée, comme celle d'un flipper.

On donne :

\begin{liste}
  \item masse volumique de l'acier : $\rho=\SI{8e3}{kg.m^{-3}}$;
  \item longueur de la rampe de lancement : $L=\SI{1,2}{m}$.
\end{liste}

Déterminer l'intervalle pour lequel la constante de raideur du ressort conviendrait.

\sol{
\textsc{Analyse stratégique de l'énoncé}


Il s'agit ici d'une résolution de problème de mécanique. Seulement deux données numériques sont fournies. Il conviendra donc, avant de se lancer dans la résolution de ce problème, de faire une analyse physique pour lister les théorèmes ou les lois à utiliser ainsi que les valeurs numériques nécessaires pour répondre à la question posée.
}
\sol{
\textsc{Rapports du jury}

\begin{liste}
  \item
  Les résolutions de problème proposées ne sont pas des questions ouvertes : une valeur chiffrée est attendue.
  
  \item
  
  Lorsqu'une donnée numérique semble manquer, il faut d'abord s'assurer qu'elle ne peut pas être mesurée / estimée à partir des documents. Si elle ne l'est pas, il est alors sans doute nécessaire de proposer un ordre de grandeur.
  
  Ce problème de mécanique de première année peut se résoudre par une méthode énergétique. Dans le cas où l'on néglige les frottements, la conservation de l'énergie mécanique (ou le théorème de la puissance mécanique) permet d'accéder à la constante de raideur du ressort.
  
  \item
  
  Il est vivement conseillé de connaître les ordres de grandeurs exigibles, les formules chimiques des espèces exigibles et plus généralement le cours en particulier de première année de CPGE de physique et de chimie qui font partie intégrante de la formation et donc de l'évaluation.
  Rapport du jury 2015
  
  Rechercher la stratégie de résolution la plus simple possible sans dénaturer le problème posé. Si le temps le permet, il sera toujours possible d'affiner le modèle choisi.
  
  Il faut lister les valeurs numériques et grandeurs physiques nécessaires mais non fournies dans l'énoncé. Pour cela on choisit des ordres de grandeurs crédibles mais qui facilitent les calculs.
\end{liste}
  }
  
  \sol{
    On choisit donc :

\begin{liste}
  \item accélération de la pesanteur à la surface de la Terre : $g=9,8 \mathrm{~m} \cdot \mathrm{s}^{-2} \approx 10 \mathrm{~m} . \mathrm{s}^{-2}$;
  \item inclinaison de la rampe de lancement du flipper : $\alpha=7^{\circ}$;
  \item intervalle d'élongation du ressort permettant de le comprimer : on estime que le ressort d'un flipper est comprimé entre $\Delta x_{\min }=6 \mathrm{~cm}$ et $\Delta x_{\max }=12 \mathrm{~cm}$;
  \item diamètre de la bille d'acier : $d=3 \mathrm{~cm}$.
\end{liste}
}

\sol{
\textsc{Rapports du jury}

Lors de la préparation : il faut commencer par repérer qualitativement les grandeurs physiques qui régissent le(s) phénomène(s) étudié(s). Dans le même temps, en sondant ses connaissances du cours, lister les expressions mathématiques des lois physiques correspondantes. Pour être productive, cette étape cruciale nécessite d'avoir compris la signification qualitative des lois physiques au programme de CPGE. Cela permet de reconnaître «l'utilité » d'une loi dans le contexte du problème posé, c'est-à-dire les liens qu'elle permettra d'établir entre les différentes grandeurs préalablement identifiées (la grandeur à déterminer, entre autres).
}

\sol{
$\hookrightarrow$ La conservation de l'énergie mécanique pour $\Delta x_{\min }$ et $\Delta x_{\max }$ permet d'obtenir I'intervalle acceptable pour la constante de raideur du ressort du flipper. Pour améliorer le modèle, on peut ensuite tenir compte d'éventuels frottements solides et/ou fluides.
\textsc{Corrigé}


En tout premier lieu, il convient de s'approprier le problème à résoudre :

On étudie le système $\{$bille en acier\}, dans le référentiel galiléen associé à la Terre. Cette bille en acier de masse $m$ sur une glissière de flipper inclinée d'un angle $\alpha$ avec l'horizontal. Le mouvement se fait sans frottement. La masse est en contact avec le ressort sans masse, de longueur à vide $x_{0}$ et de constante de raideur $k$, fixé en $O$.

Il est nécessaire de représenter sur un schéma la situation décrite dans l'énoncé.

%%%%%%%%%%%%%%%%%%%%%%%%%%%%%%%%%%%%%%%%%%%
\begin{fig}[label]
\centering
\includegraphics[width=0.5\columnwidth]{2024_05_15_889b7b71fb18f9258c1bg-2}
% \capt{}
\end{fig}
%%%%%%%%%%%%%%%%%%%%%%%%%%%%%%%%%%%%%%%
}

\sol{
\textsc{Rapports du jury}

Faire un schéma est indispensable, retenir et noter au tableau les informations nécessaires, introduire et noter au tableau les grandeurs pertinentes à la résolution.
}
\sol{
On considère que tous les frottements (solides et fluides) sont négligés. La bille en acier est donc soumise uniquement à des forces conservatives. On peut donc appliquer le théorème de la puissance mécanique :

$$
\frac{d E_{m}}{d t}=\mc{P}_{n c}=0
$$

On peut donc en déduire que l'enérgie mécanique se conserve pendant tout le mouvement de la bille. L'énergie mécanique de la bille en acier s'exprime ainsi :

$$
E_{m}=E_{c}+E_{p}=E_{c}+E_{p, p e s}+E_{p, e l a}=\frac{1}{2} m v^{2}+m g z+\text { Cte }+\frac{1}{2} k\left(x-x_{0}\right)^{2}
$$

La conservation de l'énergie mécanique va donc s'appliquer entre deux positions de la bille. On définit :

\begin{liste}
  \item L'état initial $(i)(z=0, v(t=0)=0)$ qui correspond au cas où la bille est au contact du ressort comprimé juste avant le lancement :

$$
E_{c}(i)=0 ; E_{p, p e s}(i)=0 ; E_{p, e l a}(i)=\frac{1}{2} k(\Delta x)^{2}
$$

  \item L'état final $(f)(z=h, v=0)$ qui correspond au cas où la bille est en haut de la rampe de lancement à vitesse nulle juste avant de redescendre dans le flipper : $E_{c}(f)=0 ; E_{p, p e s}(f)=m g h ; E_{p, e l a}(f)=0$.
\end{liste}
}
\sol{
\textsc{Rapports du jury}
Il faut exposer clairement la démarche envisagée pour répondre à la question posée. Le candidat doit être force de proposition et ne pas attendre que l'examinateur lui propose des pistes. Il faut cependant rester modeste et commencer par proposer des modélisations simples qui vont amener à un résultat chiffré, contrairement aux modélisations prenant en compte trop de facteurs et rendant la résolution impossible au tableau. Il faut préciser et écrire explicitement les hypothèses faites.
}

\sol{
On applique la conservation de l'énergie mécanique entre l'état initial (i) et l'état final $(f)$ :

$$
E_{m}(i)=E_{m}(f) \Leftrightarrow \frac{1}{2} k(\Delta x)^{2}=m g h
$$

On peut alors exprimer la hauteur maximale atteinte $h$. D'après le schéma, on a :

$$
h=L \sin \alpha
$$

On peut également exprimer la masse $m$ de la bille en acier. L'énoncé donne la valeur de la masse volumique $\rho$ de l'acier. En considérant une bille sphérique de rayon $r=\frac{d}{2}$, on obtient :

$$
m=\rho V_{\text {bille }}=\rho \frac{4}{3} \pi r^{3}
$$

On peut alors déterminer la constante de raideur $k$ du ressort :

$$
\frac{1}{2} k(\Delta x)^{2}=\rho \frac{4}{3} \pi r^{3} g L \sin \alpha \Leftrightarrow k=\frac{8 \rho \pi r^{3} g L \sin \alpha}{3(\Delta x)^{2}}
$$

On peut alors exprimer l'intervalle dans lequel se situe la constante de raideur $k$ :

$$
k_{\min }=\frac{8 \rho \pi r^{3} g L \sin \alpha}{3\left(\Delta x_{\max }\right)^{2}} ; k_{\max }=\frac{8 \rho \pi r^{3} g L \sin \alpha}{3\left(\Delta x_{\min }\right)^{2}} .
$$
}

\sol{
\textsc{Rapports du jury}

Il faut mettre en équations le problème. Trop souvent les candidats disent ce qu'ils pourraient faire mais ne le font pas. L'examinateur note ce qui est fait et non ce qui aurait pu être fait. Utiliser les schémas faits pour représenter le système étudié, les forces appliquées... Appliquer les lois physiques dans le cadre des hypothèses. Faire des applications numériques pour quantifier le ou les résultats et ne pas hésiter à introduire les grandeurs numériques nécessaires à la résolution.
}
\sol{
Applications numériques :

$$
\begin{gathered}
k_{\min }=\frac{8 \rho \pi r^{3} g L \sin \alpha}{3\left(\Delta x_{\max }\right)^{2}} \approx \frac{8 \times 8.10^{3} \times \not \partial \times\left(\frac{3}{2} \cdot 10^{-2}\right)^{3} \times 10 \times 1,2 \times \sin 30}{\not \supset \times\left(12 \cdot 10^{-2}\right)^{2}} \\
k_{\min } \approx \frac{64.10^{4} \times \frac{27.10^{-6}}{8} \times 1,2 \times 0,12}{1,2^{2} \cdot 10^{-2}} \approx 22 \mathrm{~N} \cdot \mathrm{m}^{-1} . \\
k_{\max }=\frac{8 \rho \pi r^{3} g L \sin \alpha}{3\left(\Delta x_{\min }\right)^{2}} \approx \frac{8 \times 8.10^{3} \times \not \partial \times\left(\frac{3}{2} \cdot 10^{-2}\right)^{3} \times 10 \times 1,2 \times \sin 30}{\not 3 \times\left(6.10^{-2}\right)^{2}} \\
k_{\max } \approx \frac{64.10^{4} \times \frac{27.10^{-6}}{8} \times 1,2 \times 0,12}{36.10^{-4}} \approx 86 \mathrm{~N} . \mathrm{m}^{-1} . \\
k_{\min }=22 \mathrm{~N} \cdot \mathrm{m}^{-1} ; k_{\max }=86 \mathrm{~N} \cdot \mathrm{m}^{-1} .
\end{gathered}
$$
}

\sol{
Amélioration du modèle :

Pour améliorer le modèle, on peut aussi tenir compte de frottements. Deux types de frottements peuvent être modélisés :

\begin{liste}
  \item Frottements fluides modélisant les frottements de l'air sur la bille en acier sphérique. On introduit alors un coefficient de frottement fluide $\lambda$. On peut alors considérer que la bille est soumise à une force de frottement $\Vec{f}$ telle que:

$$
\Vec{f}=-\lambda \Vec{v}
$$

\item Frottements solides modélisant les frottements de la bille sur la rampe de lancement. On considère alors que la bille est soumise à une réaction tangentielle $\Vec{R_{N}}$. La norme de $\Vec{R_{N}}$ s'exprime alors par la loi de Coulomb :

$$
R_{T}=\mu R_{N},
$$

$\mu$ est le coefficient de frottement solide.
\end{liste}

En estimant numériquement les différents coefficients de frottements et en tenant compte de la puissance de ces forces non conservatives dans le théorème de la puissance mécanique, on peut alors calculer les nouvelles valeurs de $k_{\min }$ et $k_{\max }$. Ces nouvelles valeurs seront obligatoirement plus élevées que celles obtenues sans frottement car les frottements ralentissent la bille en acier au cours de son mouvement le long de la rampe de lancement.
}
\sol{
\textsc{Techniques à mémoriser}

\begin{liste}
  \item[
  $\checkmark$] Il faut se souvenir qu'il est obligatoire de modéliser un problème de mécanique par un schéma clair et bien annoté et de bien définir le système étudié.
  
  Rapport du jury 2014 :
  
  En mécanique les systèmes sont mal définis ce qui entraîne souvent des erreurs de signe dans les bilans énergétiques.
  
  \item[$\checkmark$] Il faut se souvenir que la conservation de l'énergie mécanique s'obtient en appliquant le théorème de la puissance mécanique dans le cas où les frottements sont négligés.
  
  II faut se souvenir que l'énergie potentielle peut intervenir sous deux formes : énergie potentielle de pesanteur et énergie potentielle élastique.
  
  \item[$\checkmark$] Il faut se souvenir qu'une résolution de problème impose de répondre littéralement à la question posée mais aussi de quantifier le résultat proposé.
  
  Rapport du jury 2017 :
  
  Ne pas mélanger le calcul littéral et les applications numériques.
\end{liste}
}
\sol{
\textsc{Formulaire}
\begin{liste}
  \item Théorème de la puissance mécanique :
$$
\frac{d E_{m}}{d t}=\mc{P}_{n c}
$$

  \item Conservation de l'énergie mécanique :
Dans le cas où les frottements sont négligés, entre deux états $(i)$ et $(f)$ :

$$
E_{m}(i)=E_{m}(f)
$$

  \item Forces conservatives, énergie potentielle :

On dit qu'une force $\Vec{F}$ est une force conservative si elle dérive d'une énergie potentielle $E_{p}$ avec :

$$
d E_{p}=-\delta W \Longleftrightarrow E_{p}=-\int \Vec{F} . \Vec{d r} \quad \text { soit } \quad F_{x}=-\frac{d E_{p}}{d x} \text { si un seul paramètre } x
$$

  \item Énergie potentielle élastique :

$$
E_{p, \text { élastique }}=\frac{1}{2} k\left(\ell-\ell_{0}\right)^{2}=\frac{1}{2} k x^{2}
$$

  \item Énergie potentielle de pesanteur :

$$
E_{p, \text { pesanteur }}= \pm m g z+C t e
$$

Le signe + est appliqué si l'axe $z$ est vertical ascendant, le signe - est appliqué si l'axe $z$ est vertical descendant.

  \item Énergie cinétique :
  
  $$
  E_{c}=\frac{1}{2} m v^{2}
  $$
\end{liste}
}

\end{exo}
%%%%%%%%%%%%%%%%%%%%%%%%%%%%%%%%%%%

